\documentclass[12pt,reqno]{amsart}
\usepackage[utf8]{inputenc}
\usepackage{amsmath,amssymb,amsthm,mathrsfs}
\usepackage{hyperref}
\usepackage{bookmark}
\hypersetup{hidelinks,pdfstartview={FitH}}
\addtolength{\hoffset}{-1.5cm}
\addtolength{\textwidth}{3cm}
\addtolength{\voffset}{-1cm}
\addtolength{\textheight}{1.2cm}
\theoremstyle{plain}
\newtheorem{theorem}{Theorem}[section]
\newtheorem{lemma}[theorem]{Lemma}
\newtheorem{proposition}[theorem]{Proposition}
\newtheorem{corollary}[theorem]{Corollary}
\newtheorem{exttheorem}[theorem]{External Theorem}
\theoremstyle{definition}
\newtheorem{definition}[theorem]{Definition}
\theoremstyle{remark}
\newtheorem{remark}[theorem]{Remark}
\numberwithin{equation}{section}
\newcommand{\R}{\mathbb R}
\newcommand{\C}{\mathbb C}
\newcommand{\N}{\mathbb N}
\newcommand{\Z}{\mathbb Z}
\newcommand{\Sp}{\mathcal S}
\newcommand{\T}{\mathcal T}
\newcommand{\g}{\mathsf g}
\newcommand{\h}{\mathsf h}
\newcommand{\B}{\mathsf B}
\newcommand{\F}{\mathcal F}
\newcommand{\dd}{\,d}
\newcommand{\one}{\mathbf 1}
\newcommand{\br}[1]{\langle #1\rangle}
\DeclareMathOperator{\supp}{supp}
\DeclareMathOperator{\ch}{ch}
\title[Anisotropic twisted paraproduct]{The three-dimensional anisotropic twisted paraproduct}
\date{\today}
\begin{document}
\maketitle
\tableofcontents

\section{Conventions and the main statement}

Let $E_d=\R^d$ with its Euclidean norm and Lebesgue measure.
Coordinates of $E_3$ are indexed by $1,2,3$, and the input functions
by $0,1,2,3$.
Product-coordinate identifications are measure-preserving linear
equivalences. All integrals are Lebesgue integrals; Fubini's theorem
is applied to signed or complex integrals only after absolute
integrability has been established.
For $1\le p<\infty$, $\|f\|_p=(\int |f|^p)^{1/p}$, and
$\|f\|_\infty$ is the essential supremum.
For pointwise fiber constructions from $L^p$ classes, choose Borel
representatives.

\begin{definition}[Anisotropy and dilations]\label{def:anisotropy}
Fix $\alpha=(\alpha_1,\alpha_2,\alpha_3)\in\N^3$ with $\alpha_i\ge1$.
Set
\[
 A=\alpha_1+\alpha_2+\alpha_3,\qquad
 D_t x=(t^{\alpha_1}x_1,t^{\alpha_2}x_2,t^{\alpha_3}x_3),\qquad
 \rho_\alpha(\xi)=\sum_{i=1}^3|\xi_i|^{1/\alpha_i}.
\]
Here $t>0$, $\det D_t=t^A$, and
$\rho_\alpha(D_t\xi)=t\rho_\alpha(\xi)$.
The inverse and composition laws are
\[
 D_t^{-1}x=(t^{-\alpha_1}x_1,t^{-\alpha_2}x_2,t^{-\alpha_3}x_3),
 \qquad D_sD_t=D_{st},\qquad D_1=\operatorname{id}.
\]
For a one-dimensional function $v$, put
\[
 v_{(s)}(x)=s^{-1}v(x/s),\quad
 v_{s,a}(x)=v_{(s)}(x-a),\quad
 (f*_i v)(x)=\int_\R f(x-ye_i)v(y)\dd y.
\]
Also write $v_s=v_{(s)}$. Thus $(v_{1,r})_{s,a}=v_{s,a+sr}$.
\end{definition}

The Fourier transform is
$\widehat f(\xi)=\int f(x)e^{-2\pi i x\cdot\xi}\dd x$ and the inverse
transform has a positive sign in the exponential. Distributional
pairings are complex bilinear, with
$\widehat T(\phi)=T(\widehat\phi)$.
Define the decaying bracket by
\[
 \br{x}=(1+|x|)^{-1},\qquad
 \br{x}_{(s)}^N=s^{-1}(1+|x|/s)^{-N}.
\]
For translation parameters put
\[
 U(u)=1+|u_1|+|u_2|+|u_3|,\qquad R(r)=1+|r|.
\]
Unless additional dependence is displayed, constants depend only on
$\alpha$ and the fixed bumps defined below. Constants in exponent
estimates may also depend on the exponents. They never depend on a
tree, its depth, a scale truncation, $u$, or a measurable coefficient
of modulus at most one.

\begin{definition}[Multiplier and quadrilinear form]\label{def:multiplier}
Let $M\ge0$ and let $m:\R^3\to\C$ be measurable, with $m(0)=0$,
smooth off the origin, and, for every multi-index $\beta\in\N^3$ with
$|\beta|\le110$ and every $\xi\ne0$, satisfying
\begin{equation}\label{eq:symbol}
 |\partial^\beta m(\xi)|\le M\rho_\alpha(\xi)^{-\alpha\cdot\beta}.
\end{equation}
For $f_j\in\Sp(E_3;\C)$ define
\[
 P_f(t)=\int_{E_3}f_0(x)\prod_{j=1}^3f_j(x+t_je_j)\dd x,
 \qquad
 \Lambda_m(f)=\int_{E_3}m(\xi)\widehat {P_f}(-\xi)\dd\xi.
\]
Equivalently, with $K=\F^{-1}m$ as a tempered distribution,
$\Lambda_m(f)=K(P_f)$.
\end{definition}

\begin{lemma}[Well-defined multiplier pairing]\label{lem:pairing}
The function $P_f$ is Schwartz, the defining integral for $\Lambda_m$
is absolutely convergent, and the stated distributional identity holds.
Writing $\zeta=(\zeta^{(1)},\zeta^{(2)},\zeta^{(3)})\in E_9$ and
$\eta(\zeta)=(\zeta^{(1)}_1,\zeta^{(2)}_2,\zeta^{(3)}_3)$, one has
the absolutely convergent frequency representation
\begin{equation}\label{eq:frequency_form}
 \Lambda_m(f)=\int_{E_9}m(-\eta(\zeta))
   \widehat f_0\left(-\sum_{j=1}^3\zeta^{(j)}\right)
                 \prod_{j=1}^3\widehat f_j(\zeta^{(j)})\dd\zeta.
\end{equation}
\end{lemma}
\begin{proof}
The linear map
$(x,t)\mapsto(x,x+t_1e_1,x+t_2e_2,x+t_3e_3)$ is injective.
The product of the four Schwartz functions, pulled back along this map,
is Schwartz on $E_6$. Integrating in $x$ preserves the Schwartz space:
differentiate under the integral and absorb an integrable factor
$(1+|x|)^{-4}$ in every Schwartz seminorm estimate.
The case $\beta=0$ of \eqref{eq:symbol} gives $|m|\le M$.
Fourier inversion on Schwartz functions and
$\F^{-1}P_f(\xi)=\widehat {P_f}(-\xi)$ give the last assertion.
For \eqref{eq:frequency_form}, apply Fourier inversion to $f_1,f_2,f_3$
in the definition of $P_f$ and integrate $f_0$ in $x$. Fubini is
justified by $\|f_0\|_1\prod_{j=1}^3\|\widehat f_j\|_1$.
The resulting integrand has oscillatory factor
$e^{2\pi i t\cdot\eta(\zeta)}$. Separate the three diagonal
coordinates $\eta$ from the other six coordinates by a coordinate
permutation, and integrate those other six coordinates. The resulting
Schwartz function of $\eta$ is $\widehat {P_f}(\eta)$ by Fourier
inversion. Substitute it into the definition of $\Lambda_m$ and use
$\xi=-\eta$. Absolute convergence follows from $|m|\le M$ and
$\|\widehat f_0\|_\infty\prod_{j=1}^3\|\widehat f_j\|_1$.
\end{proof}

\begin{theorem}[Anisotropic paraproduct]\label{thm:main}
Suppose $4<p_0<\infty$, $1<p_1,p_2,p_3<4$, and
$\sum_{j=0}^3p_j^{-1}=1$. There exists $C_{\alpha,p}>0$ such that
every multiplier in Definition~\ref{def:multiplier} and every complex
Schwartz tuple satisfy
\[
 |\Lambda_m(f)|\le C_{\alpha,p}M\prod_{j=0}^3\|f_j\|_{p_j}.
\]
\end{theorem}

\begin{remark}[Distributional multiplier hypothesis]\label{rem:distributional}
The assumption on $K$ is that $\widehat K$ is the regular distribution
defined by the bounded function $m$, including across the origin.
Merely identifying $\widehat K$ with $m$ on $\R^3\setminus\{0\}$
would allow an extra distribution supported at the origin. For example,
$\widehat K=\delta_0$ would have zero punctured symbol and $K=1$;
its form scales like $L^6$ on inputs $f_j(x/L)$, whereas the claimed
norm bound scales like $L^3$. The regular-distribution interpretation
is therefore necessary. Setting $m(0)=0$ changes no integrals.
The proof uses multiplier derivative bounds only through order $110$.
\end{remark}

\section{Function spaces and fixed bumps}

\begin{definition}[Schwartz test functions]\label{def:schwartz}
Let $\mathbb K$ be $\R$ or $\C$. The space $\Sp(E_d;\mathbb K)$
consists of smooth functions $f:E_d\to\mathbb K$ for which
\[
 q_{N,\gamma}(f)=\sup_{x\in E_d}(1+|x|)^N|\partial^\gamma f(x)|<\infty
 \qquad(N\in\N,\ \gamma\in\N^d).
\]
These seminorms define the Schwartz topology.
Schwartz functions are smooth, so restriction to a specified
fiber is meaningful before passing to almost-everywhere classes.
Write $C_b(E_d;\mathbb K)$ for bounded continuous functions, with
the uniform norm.
\end{definition}

\begin{lemma}[Schwartz-space toolkit]\label{lem:schwartz}
The following properties hold over $\mathbb K\in\{\R,\C\}$.
\begin{enumerate}
\item $\Sp(E_d)\subset C_b(E_d)\cap L^p(E_d)$
for every $1\le p\le\infty$.
\item The Schwartz class is closed under differentiation, products,
finite tensor products, and translations. The Fourier transform and
its inverse are continuous maps of the complex Schwartz space to itself.
\item If $L:E_k\to E_n$ is injective linear, $a\in E_n$, and
$f\in\Sp(E_n)$, then $x\mapsto f(Lx+a)$ belongs to $\Sp(E_k)$.
More generally, if $f_j\in\Sp(E_{d_j})$ and the joint linear map
$x\mapsto(L_jx)_j$ is injective, then
$x\mapsto\prod_j f_j(L_jx+a_j)$ is Schwartz.
\item If $F\in\Sp(E_m\times E_l)$, then every fixed fiber
$v\mapsto F(u,v)$ is Schwartz and
\[
 g(u)=\int_{E_l}F(u,v)\dd v
 \quad\text{belongs to }\Sp(E_m),\qquad
 \partial^\beta g(u)=\int_{E_l}\partial_u^\beta F(u,v)\dd v.
\]
These operations have bounds in the corresponding Schwartz seminorms;
fixed-fiber bounds may depend on the fixed value of $u$.
In particular,
\[
 \sup_u\int_{E_l}|F(u,v)|\dd v<\infty.
\]
\item If $f\in\Sp(E_d)$ and $v\in\Sp(\R)$, then $f*_i v$
is Schwartz. Ordinary convolution of two Schwartz functions is
Schwartz as well. Schwartz functions are dense in $L^p(E_d)$ for
every finite $1\le p<\infty$.
\end{enumerate}
\end{lemma}
\begin{proof}
For the embeddings use $|f(x)|\le q_{N,0}(f)(1+|x|)^{-N}$ with
$N>d$, together with boundedness. Leibniz and the chain rule give
the derivative, product, tensor and translation assertions.
For example, the product bound is
\[
 q_{N,\gamma}(fg)\le
 \sum_{\beta\le\gamma}\binom{\gamma}{\beta}
        q_{N,\beta}(f)q_{0,\gamma-\beta}(g).
\]
For the Fourier transform, differentiation under the integral and
integration by parts give
\[
 \partial_\xi^\gamma\widehat f
       =\F\big((-2\pi ix)^\gamma f\big),\qquad
 (2\pi i\xi)^\beta\partial_\xi^\gamma\widehat f
       =\F\big(\partial_x^\beta[(-2\pi ix)^\gamma f]\big).
\]
The absolute value of the last expression is bounded by the $L^1$
norm of its argument. Bound that norm by a finite sum of Schwartz
seminorms, using the integrable weight $(1+|x|)^{-d-1}$.
Expanding $(1+\sum_i|\xi_i|)^N$ then bounds every
$q_{N,\gamma}(\widehat f)$ by a finite sum of seminorms of $f$.
This proves continuity, and changing the Fourier sign proves the
same assertion for the inverse transform.

For injective $L$, finite-dimensional norm comparison gives
$|x|\le C_L|Lx|$, hence
$1+|x|\le C_{L,a}(1+|Lx+a|)$. Apply this to each derivative in
the chain rule. For the joint-map statement first take the tensor
product of the $f_j$, and then apply the injective-pullback statement.
An individual coordinate projection need not be injective; no
Schwartz pullback claim is made for such a projection by itself.

For fiber integration, rapid decay gives, for every $N$ and $\beta$,
\[
 |\partial_u^\beta F(u,v)|
       \le q_{N+l+1,(\beta,0)}(F)
                     (1+|u|)^{-N}(1+|v|)^{-l-1}.
\]
The last factor is integrable. Dominated differentiation proves the
displayed derivative identity, and integration of this majorant gives
\[
 q_{N,\beta}(g)\le C_lq_{N+l+1,(\beta,0)}(F),\qquad
 C_l=\int_{E_l}(1+|v|)^{-l-1}\dd v<\infty.
\]
Restriction to a fixed fiber follows from injective affine pullback.
The same bound with $N=0$ gives the uniform absolute fiber integral.

For coordinate convolution, $(x,y)\mapsto(x-ye_i,y)$ is a linear
isomorphism. Thus $f(x-ye_i)v(y)$ is Schwartz jointly in $(x,y)$;
integrating in $y$ proves the assertion. Use
$(x,y)\mapsto(x-y,y)$ for ordinary convolution.
Finally, truncation followed by convolution with a smooth compactly
supported approximate identity gives density in finite $L^p$.
\end{proof}

\begin{remark}[Domains of the auxiliary forms]
The global model estimates are first proved for Schwartz inputs and
extended to finite $L^p$ spaces by density. Local cube forms are defined
on $C_b(E_3)^8$, which contains the Schwartz functions and the constant
function $1$ used in the edge slots.
All derivatives in the local telescoping identity fall on the kernels,
so boundedness and continuity of the local inputs suffice.

Absolute values and noninteger powers of Schwartz functions need not
be smooth. Expressions such as $|f|^s$ and $|\varphi|$ below are
used as measurable functions with the decay bounds inherited from
the original Schwartz functions. The bracket kernels $\theta$ and
$\Theta_t$ are positive integrable comparison weights; Schwartz
regularity is not required of them. The measurable good and bad
functions in the fiberwise decomposition are handled in $L^p$.
\end{remark}

\begin{definition}[Gaussian and bracket kernels]\label{def:gaussian}
Set
\[
 \g(x)=e^{-\pi x^2},\quad \h(x)=\g'(x)=-2\pi x\g(x),\quad
 \theta(x)=\br{x}^{10},\quad
 \Theta_t=\bigotimes_{i=1}^3\theta_{(t^{\alpha_i})}.
\]
Notice that $\int\g=1$, $\int\h=0$, $\int\theta=2/9\le1$,
$\widehat\g=\g$, and $\widehat\h(\xi)=2\pi i\xi\g(\xi)=-i\h(\xi)$.
\end{definition}

\begin{definition}[Standard bump and cone cutoffs]\label{def:bumps}
Define the standard bump $\B$ and its Fourier transform $b$ by
\[
 \B_l=\widehat{\one_{[-3/4,3/4]}}
       \prod_{i=0}^{l-1}2^{i+2}
          \widehat{\one_{[-2^{-i-3},2^{-i-3}]}},\qquad
 \B=\lim_{l\to\infty}\B_l,\qquad b=\widehat\B.
\]
Define
\[
 \Psi(\eta)=b(4(\eta-3/2))+b(4(\eta+3/2)),\qquad
 c_\Psi=\int_0^\infty\Psi(v)\h(v)^2\frac{\dd v}{v},
\]
\[
 \Phi(\eta)=\int_{|\eta|}^\infty\Psi(v)\h(v)^2\frac{\dd v}{v},
 \qquad \psi=\F^{-1}\Psi,\quad\varphi=\F^{-1}\Phi.
\]
At zero the integrand $\Psi(v)\h(v)^2/v$ is defined to be zero.
\end{definition}

\begin{lemma}[Bump properties]\label{lem:bumps}
The standard bump is real, even, and Schwartz; $0\le b\le1$,
$\supp b\subset[-1,1]$, and $b=1$ on $[-1/2,1/2]$.
Moreover $c_\Psi>0$, $\Psi$ and $\Phi$ are real, even and smooth,
\[
 \supp\Psi\subset[-2,-1]\cup[1,2],\quad
 \supp\Phi\subset[-2,2],\quad \Phi=c_\Psi\text{ on }[-1,1].
\]
The functions $\psi,\varphi$ are real even Schwartz functions.
For each fixed integer $N\ge0$, the seminorms
$\sup_x(1+|x|)^N|v^{(a)}(x)|$, $a=0,1$, are finite for
$v=\g,\h,\psi,\varphi,\g*\varphi,\h*\h*\psi$.
\end{lemma}
\begin{proof}
Put $a_i=2^{-i-3}$ and
$\mu_i=2^{i+2}\one_{[-a_i,a_i]}$. These are nonnegative
probability densities. For $l\ge1$, Fourier transformation of the finite
product, using evenness to remove reflection, gives
\[
 b_l=\widehat\B_l
     =\one_{[-3/4,3/4]}*\mu_0*\cdots*\mu_{l-1},\qquad
 \sum_{i=0}^{l-1}a_i=\frac14-2^{-l-2}.
\]
In particular $0\le b_l\le1$, its support is contained in
$[-1+2^{-l-2},1-2^{-l-2}]$, and
$b_l=1$ on $[-1/2-2^{-l-2},1/2+2^{-l-2}]$.
For fixed $\xi$, the factors
\[
 \widehat\mu_i(\xi)=\frac{\sin(2\pi a_i\xi)}{2\pi a_i\xi},
 \qquad \widehat\mu_i(0)=1,
 \qquad |1-\widehat\mu_i(\xi)|\le C a_i^2|\xi|^2
\]
show convergence of the product, uniformly on compact $\xi$ sets.
For $l\ge1$, the initial interval factor and $\widehat\mu_0$
give $|\B_l(\xi)|\le C(1+|\xi|)^{-2}$, since every remaining
factor has modulus at most one. Dominated convergence gives
$\B_l\to\B$ in $L^1$ and therefore $b_l\to b$ uniformly.
The positivity, support, and plateau properties pass to $b$.

Distributionally, the finite signed measure
\[
 \mu_i'=2^{i+2}(\delta_{-a_i}-\delta_{a_i})
 \quad\hbox{has total variation }2^{i+3}.
\]
For $l\ge\max(1,r)$, differentiate distinct factors
$\mu_0,\ldots,\mu_{r-1}$.
Convolution with the other probability densities is an $L^\infty$
contraction, whence
\[
 \|b_l^{(r)}\|_\infty
       \le\prod_{i=0}^{r-1}2^{i+3}
       =2^{r(r-1)/2+3r}.
\]
Uniform convergence transfers these weak derivative bounds to $b$.
Applying the one-dimensional fundamental theorem for weak
derivatives successively, for arbitrarily large $r$, gives a smooth
representative. It is the continuous function $b$ already obtained.
Thus $b\in C_c^\infty$, and Fourier inversion and
Lemma~\ref{lem:schwartz} give $\B\in\Sp(\R)$.
Our two translates of $b$ have supports in
$[5/4,7/4]$ and $[-7/4,-5/4]$. Positivity on a subinterval implies
$c_\Psi>0$. In addition to the defining formula, for $\eta\ne0$
the substitution $v=s|\eta|$ gives
\[
 \Phi(\eta)=\int_1^\infty\Psi(s\eta)\h(s\eta)^2\frac{\dd s}{s}.
\]
The defining formula gives $0\le\Phi\le c_\Psi$, with
$\Phi=c_\Psi$ on $[-1,1]$ and $\Phi=0$ for $|\eta|\ge2$.
For $\eta\ne0$ its derivative is
\[
 \Phi'(\eta)=-\frac{\Psi(\eta)\h(\eta)^2}{\eta}.
\]
The right side vanishes near zero, so the formula and the constant
extension prove smoothness everywhere. Inverse Fourier transformation
and Schwartz convolution give the remaining assertions.
\end{proof}

\begin{lemma}[Gaussian domination]\label{lem:domination}
For $\lambda\ge1$ and $r\in\R$,
\begin{align*}
 |\g_{\lambda,0}(v)|+|\h_{\lambda,0}(v)|&\le C\lambda^9\theta(v),\\
 |\g_{1,r}(v)|+|\h_{1,r}(v)|&\le C R(r)^{10}\theta(v).
\end{align*}
For $z\in E_4$, $u\in E_3$, one also has
\begin{equation}\label{eq:superposition}
 |\varphi(z_1-u_1)|\g(z_2)|\varphi(z_3-u_2)|\g(z_4)
 \le C U(u)^{50}\int_1^\infty
       \prod_{a=1}^4\g_{(\lambda)}(z_a)\lambda^{-47}\dd\lambda.
\end{equation}
Finally $|\psi(r\pm u_3)|\le C U(u)^{50}R(r)^{-50}$ and
$\int|\psi(r\pm u_3)|\dd r=\|\psi\|_1$.
\end{lemma}
\begin{proof}
Use $1+|v|\le\lambda(1+|v|/\lambda)$ and
$1+|v|\le R(r)(1+|v-r|)$ with order $10$ decay.
The left side of \eqref{eq:superposition} is at most
$C U(u)^{50}(1+|z|)^{-50}$. In its right side, restrict
$\lambda$ to $[1+|z|,2(1+|z|)]$; the Gaussian exponential is at
least $e^{-\pi}$ and the remaining integrand is $\lambda^{-51}$.
This proves the required lower bound for the positive superposition.
Translation invariance and order $50$ decay prove the last assertions.
\end{proof}

\section{Dyadic geometry and local sizes}

\begin{definition}[Boxes, forests, trees, and regions]\label{def:tree}
A box is encoded by $(k,n)\in\Z\times\Z^3$ and denotes
\[
 Q(k,n)=\prod_{i=1}^3[2^{k\alpha_i}n_i,2^{k\alpha_i}(n_i+1)).
\]
Set $\ell(Q)=2^k$, $|Q|=2^{kA}$, and let $c(Q)$ be its center.
Write
\[
 I_i(Q)=[a_i(Q),b_i(Q)),\qquad
 a_i(Q)=2^{k\alpha_i}n_i,\quad
 b_i(Q)=2^{k\alpha_i}(n_i+1),\quad
 c(Q)_i=2^{k\alpha_i}(n_i+\tfrac12).
\]
Its parent has level $k+1$ and indices
$\lfloor n_i/2^{\alpha_i}\rfloor$ (integer floor, also for negative
indices). Its children are
\[
 \ch(Q(k,n))=
 \left\{Q(k-1,m):m_i=2^{\alpha_i}n_i+v_i,
          \quad v_i\in\{0,\ldots,2^{\alpha_i}-1\}\right\}.
\]
A finite collection is convex if it contains every intermediate box
between two of its members ordered by inclusion. A convex tree is a
nonempty finite convex collection with a largest member $Q_\T$.
Its leaves and its collections at a fixed level are
\[
 \mathcal L(\T)=\{L:L\notin\T,\ \operatorname{parent}(L)\in\T\},
 \qquad \T_k=\{Q\in\T:\ell(Q)=2^k\},\qquad
 T_k=\bigcup_{Q\in\T_k}Q.
\]
For any finite collection $\mathcal Q$ put
\[
 \Omega_{\mathcal Q}=\bigcup_{Q\in\mathcal Q}
       Q\times(\ell(Q)/2,\ell(Q)],\qquad
 d\nu(p,t)=\dd p\,\frac{\dd t}{t}\quad(t>0).
\]
Empty collections have empty region. Upper half-open scale intervals
give exact disjointness; closures are used only in local suprema.
\end{definition}

\begin{definition}[Oriented faces and adjacent boxes]\label{def:faces}
For $\varepsilon\in\{-1,1\}$ put
\[
 q_i(Q,\varepsilon)=
 \begin{cases}a_i(Q),&\varepsilon=-1,\\b_i(Q),&\varepsilon=1,
 \end{cases}
 \qquad Q_i^\varepsilon=Q(k,n+\varepsilon e_i)
 \quad\text{if }Q=Q(k,n),
\]
where $e_i$ is the $i$th coordinate unit vector in $\Z^3$.
The closed $i$-face with outward orientation $\varepsilon$ is
\[
 \partial_i^\varepsilon Q
 =\{p:p_i=q_i(Q,\varepsilon),\ 
                  p_m\in[a_m(Q),b_m(Q)]\ (m\ne i)\}.
\]
Its two-dimensional measure means the pushforward of coordinate
Lebesgue measure under insertion of $q_i(Q,\varepsilon)$ in the
$i$th coordinate, rather than three-dimensional Lebesgue measure.
Encode each surviving face by its incident box and orientation:
\[
 \mathscr S_{i,k}
 =\{(Q,\varepsilon):Q\in\T_k,\ 
             \varepsilon\in\{-1,1\},\ Q_i^\varepsilon\notin\T\}.
\]
For $S=(Q,\varepsilon)\in\mathscr S_{i,k}$ define
\[
 Q_S=Q,\quad Q_S^{\rm out}=Q_i^\varepsilon,\quad
 \varepsilon(S)=\varepsilon,\quad q_i(S)=q_i(Q,\varepsilon),
 \qquad R_S=\prod_{m\ne i}I_m(Q).
\]
The underlying closed face of $S$ is
\[
 \{p:p_i=q_i(S),\ p_{-i}\in\overline{R_S}\},\qquad
 |R_S|=2^{k(A-\alpha_i)}.
\]
The boundary of $R_S$ has two-dimensional measure zero, so face
integrals can be taken over the half-open set $R_S$. The surviving
faces are precisely the elementary $i$-faces contained in
$\partial T_k$.
\end{definition}

\begin{lemma}[Dyadic facts and boundary packing]\label{lem:geometry}
Intersecting boxes are nested. A box has $2^A$ children, which
partition it, and the leaves of a finite convex tree partition its root.
All these sets are measurable. The surviving faces satisfy
\begin{equation}\label{eq:packing}
 \sum_k2^{kA}\#\mathscr S_{i,k}\le8|Q_\T|\quad(i=1,2,3).
\end{equation}
Faces carry two-dimensional coordinate Lebesgue measure and have area
$2^{k(A-\alpha_i)}$.
\end{lemma}
\begin{proof}
For $k\le s$, every interval of length $2^{s\alpha_i}$ is the
disjoint union of its $2^{(s-k)\alpha_i}$ dyadic subintervals of
length $2^{k\alpha_i}$. Thus a scale-$2^k$ box intersecting a
scale-$2^s$ box lies inside it. Integer division gives the stated
parent and children, and multiplication of the coordinate counts gives
$\#\ch(Q)=\prod_i2^{\alpha_i}=2^A$. At any fixed level the boxes
partition $E_3$; all boxes are Borel sets.

Given $x\in Q_\T$, follow its unique descendant chain starting at
$Q_\T$. The tree is finite, so there is a first absent child $L$;
then $x\in L\in\mathcal L(\T)$. Distinct leaves are disjoint:
if $L\subsetneq L'$ were two leaves, then
$L'\subseteq\operatorname{parent}(L')$ and
$\operatorname{parent}(L)\subseteq L'$. Convexity between these
two parents, which belong to $\T$, would force $L'\in\T$.
Hence
\begin{equation}\label{eq:leaf_partition}
 Q_\T=\bigsqcup_{L\in\mathcal L(\T)}L,\qquad
 \sum_{L\in\mathcal L(\T)}|L|=|Q_\T|.
\end{equation}

Fix $i$ and a surviving face $S\in\mathscr S_{i,k}$. There is
exactly one incident scale-$2^k$ tree box, namely $Q_S$: if the
other incident box were also in the tree, the face would not survive.
First suppose $Q_S^{\rm out}\subseteq Q_\T$. Let $P$ be the first
ancestor of $Q_S^{\rm out}$ belonging to $\T$, which exists because
the chain reaches $Q_\T$. Let $L$ be the child of $P$ containing
$Q_S^{\rm out}$. By minimality of $P$,
\[
 Q_S^{\rm out}\subseteq L,\qquad L\notin\T,\qquad
 \operatorname{parent}(L)=P\in\T,
\]
so $L\in\mathcal L(\T)$. Also $Q_S\not\subseteq L$: otherwise
$Q_S\subseteq L\subseteq P$ and convexity would give $L\in\T$.
For $m\ne i$, adjacency gives
\[
 I_m(Q_S)=I_m(Q_S^{\rm out})\subseteq I_m(L).
\]
Therefore failure of $Q_S\subseteq L$ occurs only in the $i$th
coordinate. The adjacent intervals $I_i(Q_S)$ and
$I_i(Q_S^{\rm out})$ have equal length, with the second inside
$I_i(L)$ and the first outside. Since $I_i(L)$ is a union of
intervals of that length, their common endpoint is one of the two
endpoints of $I_i(L)$. The face $S$ consequently lies in an
$i$-face of $L$.

For $L\in\mathcal L(\T)$ define $\mathscr S_{i,k}(L)$ to be the
faces for which this ancestor construction produces $L$. These sets
partition the surviving faces with outside box contained in the root.
If $\ell(L)=2^s$, then $\mathscr S_{i,k}(L)$ is empty for $k>s$.
For $k\le s$, either $i$-face of $L$ is tiled by exactly
\[
 \prod_{m\ne i}2^{(s-k)\alpha_m}
 =2^{(s-k)(A-\alpha_i)}
\]
elementary scale-$2^k$ faces. Thus
\begin{align*}
 \#\mathscr S_{i,k}(L)&\le2\,2^{(s-k)(A-\alpha_i)},\\
 \sum_k2^{kA}\#\mathscr S_{i,k}(L)
 &\le 2\sum_{k\le s}2^{kA}2^{(s-k)(A-\alpha_i)}\\
 &=2\,2^{sA}\sum_{d=0}^\infty2^{-d\alpha_i}
 =\frac{2|L|}{1-2^{-\alpha_i}}\le4|L|.
\end{align*}
Summing over the leaves and using \eqref{eq:leaf_partition} bounds
the total contribution of these faces by $4|Q_\T|$.

If $Q_S^{\rm out}\not\subseteq Q_\T$, the root is a union of
scale-$2^k$ boxes and contains $Q_S$, so $S$ lies on one of its
two $i$-faces. Write $\ell(Q_\T)=2^{s_\T}$. There are at most
$2\,2^{(s_\T-k)(A-\alpha_i)}$ such faces at level $k\le s_\T$,
and their total weighted contribution is at most
\[
 2\sum_{k\le s_\T}2^{kA}2^{(s_\T-k)(A-\alpha_i)}
 =\frac{2|Q_\T|}{1-2^{-\alpha_i}}\le4|Q_\T|.
\]
The two cases are disjoint and exhaustive, proving
\eqref{eq:packing}. All sums over faces are finite; the infinite
geometric series above are only numerical upper bounds.
\end{proof}

\begin{definition}[Local size]\label{def:size}
For $s\ge1$, $f\in C_b(E_3)$ and nonempty finite $\mathcal Q$, define
\[
 M_{\mathcal Q,s}(f)=\max_{Q\in\mathcal Q}
       (|f|^s*\Theta_{\ell(Q)}(c(Q)))^{1/s}.
\]
Write $J=\{0,1\}^3$. For $x=(x_i^b)_{1\le i\le3,\,b\in\{0,1\}}$
let $\Pi_j x=(x_1^{j_1},x_2^{j_2},x_3^{j_3})$.
Thus $E_6$ is identified with two copies of each coordinate of $E_3$,
and $\dd x=\prod_{i=1}^3\prod_{b=0}^1\dd x_i^b$.
For $F\in C_b(E_3;\R)^J$, put
\[
 \mathcal C_\T(F)=\sup_{(p,t)\in\overline{\Omega_\T}}
 \int_{E_6}\prod_{j\in J}|F_j(\Pi_j x)|
       \prod_{i,b}\theta_{t^{\alpha_i},p_i}(x_i^b)\dd x.
\]
For empty collections set the size to zero. For $L^p$ inputs with
$p>s$ use the same formula with measurable representatives; the
convolutions are finite by H\"older and are independent of those
representatives.
\end{definition}

\begin{lemma}[Size comparison and zero size]\label{lem:size}
There is $C_{\alpha,s}$ such that
\[
 M_{\T,s}(f)\le
 \sup_{(p,t)\in\overline{\Omega_\T}}(|f|^s*\Theta_t(p))^{1/s}
 \le C_{\alpha,s} M_{\T,s}(f).
\]
All suprema are finite. For a nonempty tree, zero size implies
$f=0$ almost everywhere, and pointwise when $f$ is continuous.
\end{lemma}
\begin{proof}
For $p\in\overline Q$ and $\ell(Q)/2\le t\le\ell(Q)$,
compare $t^{\alpha_i}$ with $\ell(Q)^{\alpha_i}$ and use
$|p_i-c(Q)_i|\le\ell(Q)^{\alpha_i}/2$ in each bracket factor.
For example, $\Theta_t(p-y)\le
2^A(3/2)^{30}\Theta_{\ell(Q)}(c(Q)-y)$ suffices.
The closure of a finite union is the union of the closures.
Strict positivity of $\Theta_{\ell(Q)}$ proves the zero-size assertion.
\end{proof}

\begin{lemma}[Full-cube and edge Brascamp--Lieb inequalities]\label{lem:bl}
For bounded measurable nonnegative $F_j$,
\begin{equation}\label{eq:bl4}
 \int_{E_6}\prod_{j\in J}F_j(\Pi_jx)
    \prod_{i,b}\theta_{t^{\alpha_i},p_i}(x_i^b)\dd x
 \le\prod_{j\in J}(F_j^4*\Theta_t(p))^{1/4}.
\end{equation}
If only $F_{000}=F_0$ and $F_{001}=F_1$ differ from $1$, the left
side is at most
\begin{equation}\label{eq:bl2}
 (F_0^2*\Theta_t(p))^{1/2}(F_1^2*\Theta_t(p))^{1/2}.
\end{equation}
In fact the following four-variable inequality, before inserting the
two unused coordinates, holds:
\begin{align}
 &\int_{E_4} F_0(x_1,x_2,z_0)F_1(x_1,x_2,z_1)
       \theta_{t^{\alpha_1},p_1}(x_1)
       \theta_{t^{\alpha_2},p_2}(x_2)
       \prod_{b=0}^1\theta_{t^{\alpha_3},p_3}(z_b)
       \dd x_1\dd x_2\dd z_0\dd z_1\notag\\
 &\hspace{25mm}\le
 (F_0^2*\Theta_t(p))^{1/2}(F_1^2*\Theta_t(p))^{1/2}.
 \label{eq:edge_bl_four_variables}
\end{align}
\end{lemma}
\begin{proof}
Set
\[
 \dd\mu_i^b(v)=\theta_{t^{\alpha_i},p_i}(v)\dd v,
 \qquad \dd\mu(x)=\prod_{i,b}\dd\mu_i^b(x_i^b),\qquad
 m_\theta=\int_\R\theta(v)\dd v=\frac29.
\]
Let $\bar j=(1-j_1,1-j_2,1-j_3)$, $J_0=\{0\}\times\{0,1\}^2$,
and define $B_j(x)=F_j(\Pi_jx)F_{\bar j}(\Pi_{\bar j}x)$.
Every vertex belongs to exactly one of these four opposite pairs, so
the left side of \eqref{eq:bl4} equals $\int\prod_{j\in J_0}B_j\dd\mu$.
Two applications of Cauchy--Schwarz give
\begin{align*}
 \int\prod_{j\in J_0}B_j\dd\mu
 &\le\left(\int B_{000}^2B_{001}^2\dd\mu\right)^{1/2}
       \left(\int B_{010}^2B_{011}^2\dd\mu\right)^{1/2}\\
 &\le\prod_{j\in J_0}\left(\int B_j^4\dd\mu\right)^{1/4}.
\end{align*}
For each $j$, the coordinate sets used by $\Pi_j$ and $\Pi_{\bar j}$
are disjoint and together exhaust the six variables. Fubini therefore
gives the exact factorization
\begin{align*}
 \int B_j^4\dd\mu
 &=\left(\int_{E_3}F_j(y)^4
                    \prod_i\theta_{t^{\alpha_i},p_i}(y_i)\dd y\right)
   \left(\int_{E_3}F_{\bar j}(z)^4
                    \prod_i\theta_{t^{\alpha_i},p_i}(z_i)\dd z\right)\\
 &=(F_j^4*\Theta_t)(p)(F_{\bar j}^4*\Theta_t)(p).
\end{align*}
Here evenness of $\theta$ identifies the weighted integrals with
the displayed convolutions. Multiplying over $J_0$ proves
\eqref{eq:bl4}.

For the four-variable assertion set
\begin{align*}
 \dd\eta&=\theta_{t^{\alpha_1},p_1}(x_1)
          \theta_{t^{\alpha_2},p_2}(x_2)
          \prod_{b=0}^1\theta_{t^{\alpha_3},p_3}(z_b)
          \dd x_1\dd x_2\dd z_0\dd z_1,\\
 B_b&=F_b(x_1,x_2,z_b).
\end{align*}
Cauchy--Schwarz and integration in the unused $z$ variable give
\begin{align*}
 \int B_0B_1\dd\eta
 &\le\left(\int B_0^2\dd\eta\right)^{1/2}
       \left(\int B_1^2\dd\eta\right)^{1/2},\\
 \int B_0^2\dd\eta
 &=\left(\int_\R\theta_{t^{\alpha_3},p_3}(z_1)\dd z_1\right)
   \int_{E_3}F_0(x_1,x_2,z_0)^2
       \theta_{t^{\alpha_1},p_1}(x_1)\\
 &\hspace{10mm}\times
       \theta_{t^{\alpha_2},p_2}(x_2)
       \theta_{t^{\alpha_3},p_3}(z_0)\dd x_1\dd x_2\dd z_0\\
 &=m_\theta(F_0^2*\Theta_t)(p).
\end{align*}
The corresponding identity for $B_1$ has the same factor
$m_\theta\le1$, proving \eqref{eq:edge_bl_four_variables}.
Finally, in the six-variable edge form, integration in $x_1^1,x_2^1$
adds the factor $m_\theta^2\le1$. This proves \eqref{eq:bl2}.
\end{proof}

\section{Cubical telescoping with boundary terms}

\begin{definition}[Local cube forms]\label{def:cube}
For a tuple $F\in C_b(E_3;\R)^J$, a finite collection $\mathcal Q$,
$\lambda\ge1$, $r\in\R$, $t>0$, and $p\in E_3$, set
\[
 (\mu_1,s_1)=(\lambda,0),\quad(\mu_2,s_2)=(\lambda,0),\quad
 (\mu_3,s_3)=(1,r),
\]
\[
 G_i(v)=\g_{\mu_i t^{\alpha_i},p_i+s_it^{\alpha_i}}(v),\quad
 H_i(v)=\h_{\mu_i t^{\alpha_i},p_i+s_it^{\alpha_i}}(v),\quad
 \mathcal F_F(x)=\prod_{j\in J}F_j(\Pi_jx).
\]
Dependence on $t,p$ is suppressed only in $G_i,H_i$.
Define
\begin{equation}\label{eq:cubeform}
 \Theta^{(i)}_{\mathcal Q,\lambda,r}(F)=
 \int_{\Omega_{\mathcal Q}}\int_{E_6}\mathcal F_F(x)
 H_i(x_i^0)H_i(x_i^1)
 \prod_{m\ne i}G_m(x_m^0)G_m(x_m^1)\dd x\,d\nu(p,t).
\end{equation}
In coordinates,
\begin{align}
 \Theta^{(1)}_{\mathcal Q,\lambda,r}(F)
 &=\int_{\Omega_{\mathcal Q}}\int_{E_6}\mathcal F_F(x)
       \prod_{b=0}^1\h_{\lambda t^{\alpha_1},p_1}(x_1^b)
       \prod_{b=0}^1\g_{\lambda t^{\alpha_2},p_2}(x_2^b)\notag\\
 &\hspace{12mm}\times\prod_{b=0}^1\g_{t^{\alpha_3},p_3+rt^{\alpha_3}}(x_3^b)
       \dd x\,d\nu(p,t),\label{eq:cubeform_first}\\
 \Theta^{(2)}_{\mathcal Q,\lambda,r}(F)
 &=\int_{\Omega_{\mathcal Q}}\int_{E_6}\mathcal F_F(x)
       \prod_{b=0}^1\g_{\lambda t^{\alpha_1},p_1}(x_1^b)
       \prod_{b=0}^1\h_{\lambda t^{\alpha_2},p_2}(x_2^b)\notag\\
 &\hspace{12mm}\times\prod_{b=0}^1\g_{t^{\alpha_3},p_3+rt^{\alpha_3}}(x_3^b)
       \dd x\,d\nu(p,t),\label{eq:cubeform_second}\\
 \Theta^{(3)}_{\mathcal Q,\lambda,r}(F)
 &=\int_{\Omega_{\mathcal Q}}\int_{E_6}\mathcal F_F(x)
       \prod_{b=0}^1\g_{\lambda t^{\alpha_1},p_1}(x_1^b)
       \prod_{b=0}^1\g_{\lambda t^{\alpha_2},p_2}(x_2^b)\notag\\
 &\hspace{12mm}\times\prod_{b=0}^1\h_{t^{\alpha_3},p_3+rt^{\alpha_3}}(x_3^b)
       \dd x\,d\nu(p,t).\label{eq:cubeform_third}
\end{align}
Define the vertex map
\[
 \rho_i^b(j)=j[i\mapsto b]
 =(j_1,\ldots,j_{i-1},b,j_{i+1},\ldots,j_3)
\]
and the coordinate-copy operator
\[
 (\mathsf C_i^bF)_j=F_{\rho_i^b(j)}.
\]
It obeys
\[
 (\mathsf C_i^b)^2=\mathsf C_i^b,\qquad
 \mathsf C_i^a\mathsf C_q^b=\mathsf C_q^b\mathsf C_i^a
 \quad(i\ne q).
\]
These equalities follow slot by slot from the analogous identities
for replacement of distinct coordinates of a vertex.
\end{definition}

\begin{lemma}[Local integrability and differentiation]\label{lem:local_integrability}
For bounded continuous real $F_j$ and a finite collection,
\eqref{eq:cubeform} is absolutely integrable and multilinear.
The same holds with a Gaussian replaced by its first scale or center
derivative, and for the face and endpoint integrals defined below.
Differentiation under the spatial integrals is valid on compact
positive scale intervals and bounded sets of centers.
\end{lemma}
\begin{proof}
The scale lies in $[a,b]\subset(0,\infty)$ and the centers in a
bounded set. Gaussian derivatives, with these parameters fixed in a
compact set, admit one common integrable rapidly decreasing majorant
in all six variables. Multiply by $\prod_j\|F_j\|_\infty$.
The same argument works on each coordinate face; its parameter domain
has finite area. Dominated convergence proves the derivative statements.
\end{proof}

\begin{definition}[Endpoint and boundary forms]\label{def:boundary}
Write
\[
 D_i=G_i(x_i^0)G_i(x_i^1),\quad P_i=H_i(x_i^0)H_i(x_i^1),
\]
\[
 E_i=\frac{\mu_i}{2\pi}
       \big(H_i(x_i^0)G_i(x_i^1)+G_i(x_i^0)H_i(x_i^1)\big)-s_iD_i,
 \qquad K_{t,p}(x)=\prod_iD_i.
\]
For $Q=I_1\times I_2\times I_3$ write
\[
 I_{-i}=\prod_{m\ne i}I_m,\qquad
 \dd p_{-i}=\prod_{m\ne i}\dd p_m,\qquad
 [U]_{p_i=a}^{p_i=b}=U|_{p_i=b}-U|_{p_i=a}.
\]
Define the interior and boundary center integrals
\[
 U_{Q,i}(t,x)=\int_Q P_i\prod_{m\ne i}D_m\dd p,
\]
\[
 V_{Q,i}(t,x)=\int_{I_{-i}}
       [E_i]_{p_i=\inf I_i}^{p_i=\sup I_i}\prod_{m\ne i}D_m\dd p_{-i}.
\]
For $S\in\mathscr S_{i,k}$, define the face integral
\[
 \mathcal J_{i,S}(t;F)
 =\int_{R_S}\int_{E_6}\mathcal F_F(x)
       E_i\big|_{p_i=q_i(S)}\prod_{m\ne i}D_m\dd x\dd p_{-i}.
\]
Set
\[
 \Xi_{\mathcal Q}(F)=\pi\sum_{Q\in\mathcal Q}
       \int_Q\int_{E_6}\mathcal F_F(x)K_{\ell(Q),p}(x)\dd x\dd p,
\]
\[
 \mathcal B_{\mathcal Q}(F)=-\pi\sum_{Q\in\mathcal Q}\sum_i\alpha_i
       \int_{\ell(Q)/2}^{\ell(Q)}t^{\alpha_i-1}
         \int_{E_6}\mathcal F_F(x)V_{Q,i}(t,x)\dd x\dd t,
\]
\[
 \mathcal R_\T(F)=\Xi_{\mathcal L(\T)}(F)-\Xi_{\{Q_\T\}}(F)
                         +\mathcal B_\T(F).
\]
Each form also depends on $\lambda,r$.
The three boundary kernels are
\begin{align*}
 E_1&=\frac{\lambda}{2\pi}
       \big(H_1(x_1^0)G_1(x_1^1)+G_1(x_1^0)H_1(x_1^1)\big),\\
 E_2&=\frac{\lambda}{2\pi}
       \big(H_2(x_2^0)G_2(x_2^1)+G_2(x_2^0)H_2(x_2^1)\big),\\
 E_3&=\frac{1}{2\pi}
       \big(H_3(x_3^0)G_3(x_3^1)+G_3(x_3^0)H_3(x_3^1)\big)
                    -rG_3(x_3^0)G_3(x_3^1).
\end{align*}
All sums in these definitions are finite, and all forms are multilinear
in the eight real inputs on $C_b(E_3)^J$.
\end{definition}

\begin{lemma}[One-dimensional differential identity]\label{lem:oned_telescoping}
For $I=[a,b]$ with $a\le b$ and fixed $x_i^0,x_i^1$,
\begin{equation}\label{eq:oned_telescoping}
 \int_I-t\partial_tD_i\dd p_i
 =\frac{\alpha_i}{\pi}\int_IP_i\dd p_i
                  +\alpha_i t^{\alpha_i}[E_i]_{\inf I}^{\sup I}.
\end{equation}
\end{lemma}
\begin{proof}
For a Schwartz function $v$ let $v^\sharp(z)=zv(z)$. For any
$\alpha>0$, direct differentiation of
$v_{t^\alpha,p}(u)=t^{-\alpha}v(t^{-\alpha}(u-p))$ gives
\begin{align}
 -t\partial_tv_{t^\alpha,p}(u)
   &=\alpha v_{t^\alpha,p}(u)
       +\alpha((v')^\sharp)_{t^\alpha,p}(u),\label{eq:scale_derivative}\\
 ((v')^\sharp)_{t^\alpha,p}(u)
   &=-t^\alpha\partial_p(v^\sharp)_{t^\alpha,p}(u)
       -v_{t^\alpha,p}(u),\label{eq:center_derivative}\\
 \partial_pv_{t^\alpha,p}(u)&=-t^{-\alpha}(v')_{t^\alpha,p}(u).
                                                        \label{eq:center_derivative_v}
\end{align}
Substituting \eqref{eq:center_derivative} in
\eqref{eq:scale_derivative}, multiplying by
$v_{t^\alpha,p}(u_1)$, and integrating by parts on $I=[a,b]$ yields
\begin{align}
 &\int_I(-t\partial_tv_{t^\alpha,p}(u_0))v_{t^\alpha,p}(u_1)\dd p\notag\\
 &\qquad=-\alpha t^\alpha[(v^\sharp)_{t^\alpha,p}(u_0)
                       v_{t^\alpha,p}(u_1)]_{p=a}^{p=b}\notag\\
 &\qquad\quad-\alpha\int_I(v^\sharp)_{t^\alpha,p}(u_0)
                       (v')_{t^\alpha,p}(u_1)\dd p.
                                             \label{eq:scale_ibp}
\end{align}
Every factor is smooth on the compact interval, so this integration
by parts is an ordinary one-dimensional identity.

For $v=\g_{\mu,s}$ use $\h(z)=-2\pi z\g(z)$ to obtain
\[
 v'=\mu^{-1}\h_{\mu,s},\qquad
 v^\sharp=-\frac{\mu}{2\pi}\h_{\mu,s}+s\g_{\mu,s},\qquad
 H_i=-\mu_it^{\alpha_i}\partial_{p_i}G_i.
\]
Fix $i$, put $u_b=x_i^b$, and apply \eqref{eq:scale_ibp} with
$v=\g_{\mu_i,s_i}$ and $\alpha=\alpha_i$. Add the result with
$u_0,u_1$ exchanged. With
\[
 W_i=H_i(x_i^0)G_i(x_i^1)+G_i(x_i^0)H_i(x_i^1)
\]
the resulting terms, before the last cancellation, are
\begin{align}
 \int_I-t\partial_tD_i\dd p_i
 &=\frac{\alpha_i}{\pi}\int_IP_i\dd p_i
    +\frac{\alpha_i\mu_i t^{\alpha_i}}{2\pi}[W_i]_{a}^{b}\notag\\
 &\quad-2\alpha_i s_it^{\alpha_i}[D_i]_{a}^{b}
    -\frac{\alpha_i s_i}{\mu_i}\int_IW_i\dd p_i.
                                                  \label{eq:oned_before_cancel}
\end{align}
The center derivative formula implies
\[
 W_i=-\mu_it^{\alpha_i}\partial_{p_i}D_i,
 \qquad
 -\frac{\alpha_i s_i}{\mu_i}\int_IW_i\dd p_i
       =\alpha_i s_it^{\alpha_i}[D_i]_{a}^{b}.
\]
Substitution in \eqref{eq:oned_before_cancel} changes the coefficient
of $s_iD_i$ from $-2$ to $-1$, leaving exactly
\[
 \frac{\alpha_i}{\pi}\int_IP_i\dd p_i
 +\alpha_it^{\alpha_i}
        \left[\frac{\mu_i}{2\pi}W_i-s_iD_i\right]_{a}^{b}.
\]
This is \eqref{eq:oned_telescoping}.
\end{proof}

\begin{lemma}[Finite telescoping identity]\label{lem:telescoping_identity}
For a single box and then a finite convex tree,
\begin{align}
 \sum_i\alpha_i\Theta^{(i)}_{\{Q\},\lambda,r}(F)
   &=\Xi_{\ch(Q)}(F)-\Xi_{\{Q\}}(F)+\mathcal B_{\{Q\}}(F),\label{eq:single_tel}\\
 \sum_i\alpha_i\Theta^{(i)}_{\T,\lambda,r}(F)&=\mathcal R_\T(F).
                                                        \label{eq:tree_tel}
\end{align}
\end{lemma}
\begin{proof}
Fix $Q$ with $\ell(Q)=2^k$. The product rule and integration of
\eqref{eq:oned_telescoping} in the other two center coordinates give
\begin{align}
 \int_Q-t\partial_tK_{t,p}(x)\dd p
 &=\sum_{i=1}^3\int_{I_{-i}}
     \left(\int_{I_i}-t\partial_tD_i\dd p_i\right)
                       \prod_{m\ne i}D_m\dd p_{-i}\notag\\
 &=\frac1\pi\sum_{i=1}^3\alpha_iU_{Q,i}(t,x)
       +\sum_{i=1}^3\alpha_it^{\alpha_i}V_{Q,i}(t,x).
                                                    \label{eq:box_derivative}
\end{align}
The fundamental theorem of calculus yields
\begin{align}
 &\int_QK_{2^{k-1},p}(x)\dd p-\int_QK_{2^k,p}(x)\dd p\notag\\
 &\qquad=\int_{2^{k-1}}^{2^k}\int_Q-t\partial_tK_{t,p}(x)
                                      \dd p\,\frac{\dd t}{t}.
                                                    \label{eq:box_scale_endpoints}
\end{align}
The local-integrability lemma justifies differentiating under the
center integral and subsequently multiplying by $\mathcal F_F(x)$
and integrating in $x$. By definition,
\[
 \Theta^{(i)}_{\{Q\},\lambda,r}(F)
 =\int_{2^{k-1}}^{2^k}\int_{E_6}
       \mathcal F_F(x)U_{Q,i}(t,x)\dd x\,\frac{\dd t}{t}.
\]
Every child $L\in\ch(Q)$ has $\ell(L)=2^{k-1}$ and the children
partition $Q$. Therefore
\[
 \Xi_{\ch(Q)}(F)
 =\pi\int_Q\int_{E_6}\mathcal F_F(x)K_{2^{k-1},p}(x)\dd x\dd p.
\]
Multiply \eqref{eq:box_scale_endpoints} by $\pi\mathcal F_F(x)$,
integrate in $x$, and substitute \eqref{eq:box_derivative}. Moving
the boundary contribution to the other side gives
\eqref{eq:single_tel}, with exactly the negative sign in the
definition of $\mathcal B_{\{Q\}}$.

Sum \eqref{eq:single_tel} over $Q\in\T$. If
$P\in\T\setminus\{Q_\T\}$, its parent lies between $P$ and
$Q_\T$ and hence belongs to $\T$ by convexity. The negative term
$-\Xi_{\{P\}}$ in the summand indexed by $P$ cancels its unique
positive occurrence inside $\Xi_{\ch(\operatorname{parent}(P))}$.
The root has no parent in the tree and supplies the only uncancelled
negative endpoint term. The uncancelled positive terms are indexed
by children of tree boxes which do not belong to the tree, exactly
$\mathcal L(\T)$. Thus the finite reindexing identity is
\[
 \sum_{Q\in\T}\big(\Xi_{\ch(Q)}-\Xi_{\{Q\}}\big)
   =\Xi_{\mathcal L(\T)}-\Xi_{\{Q_\T\}},\qquad
 \sum_{Q\in\T}\mathcal B_{\{Q\}}=\mathcal B_\T.
\]
This proves \eqref{eq:tree_tel}.
\end{proof}

\begin{lemma}[Cancellation of interior faces]\label{lem:face_cancellation}
For every finite convex tree,
\begin{equation}\label{eq:boundary_face_expansion}
 \mathcal B_\T(F)
 =-\pi\sum_{i=1}^3\alpha_i\sum_{k\in\Z}
       \int_{2^{k-1}}^{2^k}t^{\alpha_i-1}
       \sum_{S\in\mathscr S_{i,k}}
             \varepsilon(S)\mathcal J_{i,S}(t;F)\dd t.
\end{equation}
Only levels occupied by the finite tree contribute to this sum.
\end{lemma}
\begin{proof}
Expand the endpoint brackets defining $V_{Q,i}$. An upper face has
coefficient $+1$ and a lower face coefficient $-1$. If both
$Q$ and $Q_i^+$ belong to $\T_k$, the upper face of $Q$ and lower
face of $Q_i^+$ have the same center coordinate $p_i$ and the same
transverse rectangle $R_S$. Their kernels $E_i\prod_{m\ne i}D_m$
also agree, since they depend on $p,t,x$, not on the box containing
the center. The two face integrals cancel exactly. The same argument
applies to every adjacent pair, each counted once. A face cannot have
more than two incident boxes at a fixed level. Thus the uncancelled
terms are exactly the records in $\mathscr S_{i,k}$ and
\[
 \sum_{Q\in\T_k}\int_{E_6}\mathcal F_F(x)V_{Q,i}(t,x)\dd x
       =\sum_{S\in\mathscr S_{i,k}}
                       \varepsilon(S)\mathcal J_{i,S}(t;F).
\]
The intersections of transverse face edges have measure zero and
do not alter this identity. Multiply by
$-\pi\alpha_it^{\alpha_i-1}$, integrate over the scale interval,
and sum in $i,k$ to obtain \eqref{eq:boundary_face_expansion}.
\end{proof}

\begin{lemma}[Remainder estimate]\label{lem:remainder}
There exists $C_\alpha>0$ such that, for every $\lambda\ge1$, $r$,
finite convex tree, and bounded continuous real tuple,
\begin{equation}\label{eq:remainder}
 |\mathcal R_\T(F)|\le C_\alpha\lambda^{37}R(r)^{21}
                        |Q_\T|\mathcal C_\T(F).
\end{equation}
\end{lemma}
\begin{proof}
For $\lambda\ge1$, rapid decay gives
\[
 |\g_{\lambda,0}(z)|+|\h_{\lambda,0}(z)|
 \le C\lambda^{-1}(1+|z|/\lambda)^{-10}
 \le C\lambda^9\theta(z).
\]
The elementary inequality
$1+|z|\le(1+|r|)(1+|z-r|)$ similarly gives
\[
 |\g_{1,r}(z)|+|\h_{1,r}(z)|\le CR(r)^{10}\theta(z).
\]
After dilation and translation,
\begin{align}
 |G_i(v)|+|H_i(v)|&\le C\lambda^9
                      \theta_{t^{\alpha_i},p_i}(v)\quad(i=1,2),\notag\\
 |G_3(v)|+|H_3(v)|&\le CR(r)^{10}
                      \theta_{t^{\alpha_3},p_3}(v).
                                               \label{eq:cube_kernel_domination}
\end{align}
For every choice $\Gamma_{i,b}\in\{G_i,H_i\}$, multiplication of
these six inequalities gives
\begin{equation}\label{eq:cube_integrand_control}
 \int_{E_6}|\mathcal F_F(x)|\prod_{i,b}|\Gamma_{i,b}(x_i^b)|\dd x
 \le C\lambda^{36}R(r)^{20}\mathcal C_\T(F)
\end{equation}
when $(p,t)\in\overline{\Omega_\T}$. There are four scaled factors
contributing $\lambda^9$ each and two translated factors contributing
$R(r)^{10}$ each.

For $p\in Q_\T$ the point $(p,\ell(Q_\T))$ belongs to
$\overline{\Omega_\T}$. If $L\in\mathcal L(\T)$ has parent $P$,
then $P\in\T$, $L\subseteq P$ and $\ell(L)=\ell(P)/2$, so
\[
 (p,\ell(L))\in
 \overline{P\times(\ell(P)/2,\ell(P)]}
       \subseteq\overline{\Omega_\T}\qquad(p\in L).
\]
Apply \eqref{eq:cube_integrand_control} at these parameters and
integrate in $p$. The leaf partition gives
\begin{align}
 |\Xi_{\{Q_\T\}}(F)|+|\Xi_{\mathcal L(\T)}(F)|
 &\le C\lambda^{36}R(r)^{20}\mathcal C_\T(F)
                      \left(|Q_\T|+\sum_{L\in\mathcal L(\T)}|L|\right)
                                                       \notag\\
 &\le C\lambda^{36}R(r)^{20}|Q_\T|\mathcal C_\T(F).
                                                 \label{eq:endpoint_bound}
\end{align}

It remains to estimate the face sum
\eqref{eq:boundary_face_expansion}. If $S\in\mathscr S_{i,k}$,
$p_{-i}\in R_S$, and $2^{k-1}\le t\le2^k$, insertion of
$q_i(S)$ gives a point $(p,t)\in\overline{\Omega_\T}$ because
$S$ is a face of $Q_S\in\T_k$. For $i=1,2$ the two terms of
$E_i$ each have coefficient $\lambda/(2\pi)$. For $i=3$ the
$HG,GH$ terms have coefficient $1/(2\pi)$ and the $GG$ term has
coefficient $-r$. Hence \eqref{eq:cube_integrand_control} gives
\begin{align}
 |\mathcal J_{i,S}(t;F)|
 &\le C\lambda^{37}R(r)^{20}|R_S|\mathcal C_\T(F)
                 \quad(i=1,2),\notag\\
 |\mathcal J_{3,S}(t;F)|
 &\le C\lambda^{36}R(r)^{21}|R_S|\mathcal C_\T(F).
                                                  \label{eq:face_kernel_bound}
\end{align}
In each coordinate the scale integral is exactly
\[
 \alpha_i\int_{2^{k-1}}^{2^k}t^{\alpha_i-1}\dd t
       =(1-2^{-\alpha_i})2^{k\alpha_i}\le2^{k\alpha_i}.
\]
Combining this with $|R_S|=2^{k(A-\alpha_i)}$,
$\lambda\ge1$, $R(r)\ge1$, and setting
$N_{i,k}=\#\mathscr S_{i,k}$, we obtain
\begin{align}
 |\mathcal B_\T(F)|
 &\le C\lambda^{37}R(r)^{21}\mathcal C_\T(F)
                  \sum_{i=1}^3\sum_k2^{kA}N_{i,k}\notag\\
 &\le C\lambda^{37}R(r)^{21}|Q_\T|\mathcal C_\T(F).
                                                  \label{eq:boundary_bound}
\end{align}
The last step is \eqref{eq:packing} for each coordinate. Finally,
the defining equality
$\mathcal R_\T=\Xi_{\mathcal L(\T)}-\Xi_{\{Q_\T\}}+\mathcal B_\T$
and \eqref{eq:endpoint_bound}, \eqref{eq:boundary_bound} prove
\eqref{eq:remainder}.
\end{proof}

\begin{corollary}[Two remainder norms]\label{cor:remainder_norms}
Write $D_{\lambda,r}=C_\alpha\lambda^{37}R(r)^{21}$ with
$C_\alpha$ enlarged if necessary. Then
\[
 |\mathcal R_\T(F)|\le D_{\lambda,r}|Q_\T|
                   \prod_{j\in J}M_{\T,4}(F_j).
\]
For a tuple with only slots $000,001$ different from $1$, one also has
\[
 |\mathcal R_\T(F)|\le D_{\lambda,r}|Q_\T|
                    M_{\T,2}(F_{000})M_{\T,2}(F_{001}).
\]
\end{corollary}
\begin{proof}
Apply \eqref{eq:bl4} or \eqref{eq:bl2} inside the supremum defining
$\mathcal C_\T$, and then Lemma~\ref{lem:size}.
\end{proof}

\begin{lemma}[Cubical Cauchy--Schwarz and positivity]\label{lem:cube_cs}
For real bounded continuous tuples,
\[
 |\Theta^{(i)}_\T(F)|^2\le
 \Theta^{(i)}_\T(\mathsf C_i^0F)\Theta^{(i)}_\T(\mathsf C_i^1F),
 \qquad \Theta^{(i)}_\T(\mathsf C_i^bF)\ge0.
\]
\end{lemma}
\begin{proof}
Write $x_{\widehat i}=(x_m^a)_{m\ne i,\,a\in\{0,1\}}\in E_4$
and introduce the positive measure
\[
 \dd\sigma_{i,p,t}(x_{\widehat i})
      =\prod_{m\ne i}\prod_{a=0}^1G_m(x_m^a)\dd x_{\widehat i},
 \qquad
 \dd\sigma_i=\one_{\Omega_\T}(p,t)
          \dd\sigma_{i,p,t}(x_{\widehat i})\dd p\,\frac{\dd t}{t}.
\]
Positivity follows from $\g\ge0$, even though $H_i$ changes sign.
If $j_i=b$, define $\Pi_{j,i}(x_{\widehat i},v)$ by placing $v$
in coordinate $i$ and $x_m^{j_m}$ in coordinate $m\ne i$.
For $b\in\{0,1\}$ set
\[
 A_b(x_{\widehat i},p,t)
 =\int_\R H_i(v)\prod_{j\in J:\,j_i=b}
                     F_j(\Pi_{j,i}(x_{\widehat i},v))\dd v.
\]
These functions are measurable. Since there are four factors in
each product and $\int|H_i|=\|\h\|_1$, they satisfy
\[
 |A_b|\le\|\h\|_1\prod_{j:j_i=b}\|F_j\|_\infty.
\]
Furthermore, $\int G_m=1$ and
$\nu(\Omega_\T)=\sum_{Q\in\T}|Q|\log2<\infty$, so
$A_b\in L^2(\sigma_i)$. Fubini separates the variables
$x_i^0,x_i^1$ in the original form and gives
\[
 \Theta^{(i)}_\T(F)
   =\int A_0A_1\dd\sigma_i,
 \qquad
 |\Theta^{(i)}_\T(F)|^2\le\mathcal I_0\mathcal I_1,
 \qquad \mathcal I_b=\int A_b^2\dd\sigma_i.
\]

To identify $\mathcal I_b$, expand the square in $A_b^2$ and name
its two independent integration variables $x_i^0,x_i^1$. The product
of input functions then equals
\begin{align*}
 \prod_{a=0}^1\prod_{j:j_i=b}
                  F_j(\Pi_{\rho_i^a(j)}x)
 &=\prod_{k\in J}F_{\rho_i^b(k)}(\Pi_kx)\\
 &=\prod_{k\in J}(\mathsf C_i^bF)_k(\Pi_kx).
\end{align*}
The first equality is the bijection
$(a,j)\mapsto k=\rho_i^a(j)$ from
$\{0,1\}\times\{j:j_i=b\}$ to $J$, with inverse
$k\mapsto(k_i,\rho_i^b(k))$. The other factors are exactly
$H_i(x_i^0)H_i(x_i^1)\prod_{m\ne i,a}G_m(x_m^a)$.
Consequently
\[
 \mathcal I_b=\Theta^{(i)}_\T(\mathsf C_i^bF)\ge0.
\]
Substitution proves both conclusions.
\end{proof}

\begin{proposition}[Full cube tree estimate]\label{prop:cube_tree}
For $i=1,2,3$,
\begin{equation}\label{eq:cube_tree}
 |\Theta^{(i)}_{\T,\lambda,r}(F)|\le
 C_\alpha\lambda^{37}R(r)^{21}|Q_\T|
                           \prod_{j\in J}M_{\T,4}(F_j).
\end{equation}
\end{proposition}
\begin{proof}
If $M_{\T,4}(F_j)=0$ for some $j$, Lemma~\ref{lem:size} implies
$F_j=0$ and multilinearity gives the estimate. Otherwise put
\[
 m_j=M_{\T,4}(F_j)>0,\qquad \widetilde F_j=m_j^{-1}F_j.
\]
Then $M_{\T,4}(\widetilde F_j)=1$ and
\[
 \Theta^{(i)}_\T(F)
     =\left(\prod_{j\in J}m_j\right)
                     \Theta^{(i)}_\T(\widetilde F).
\]
It suffices to prove the estimate for the normalized tuple, which we
now denote by $F$. Each slot of every copied tuple is one of these
normalized functions, so its size is also one. Thus
Corollary~\ref{cor:remainder_norms} gives
\[
 |\mathcal R_\T(\mathsf C_{i_n}^{b_n}\cdots
                     \mathsf C_{i_1}^{b_1}F)|\le D,
 \qquad D=D_{\lambda,r}|Q_\T|,
\]
for every finite sequence of copies. We suppress the fixed parameters
$\T,\lambda,r$ in the following displayed calculations.

Fix $i$. The first Cauchy--Schwarz step is
\begin{equation}\label{eq:cube_first_cs}
 |\Theta^{(i)}(F)|^2
 \le\Theta^{(i)}(F^0)\Theta^{(i)}(F^1),\qquad
 F^a=\mathsf C_i^aF\quad(a=0,1).
\end{equation}
For each $a$, positivity and telescoping give
\begin{align}
 0\le\alpha_i\Theta^{(i)}(F^a)
 &=\mathcal R(F^a)-\sum_{q\ne i}\alpha_q\Theta^{(q)}(F^a)\notag\\
 &\le D+\sum_{q\ne i}\alpha_q|\Theta^{(q)}(F^a)|.
                                                  \label{eq:cube_first_reduction}
\end{align}
Fix one of the two coordinates $q\ne i$ and let $p$ be the unique
element of $\{1,2,3\}\setminus\{i,q\}$. Set
\[
 F^{a,b}=\mathsf C_q^bF^a
        =\mathsf C_q^b\mathsf C_i^aF\quad(b=0,1).
\]
The second Cauchy--Schwarz step is
\begin{equation}\label{eq:cube_second_cs}
 |\Theta^{(q)}(F^a)|^2
       \le\Theta^{(q)}(F^{a,0})\Theta^{(q)}(F^{a,1}).
\end{equation}
Commutativity and idempotence of copies give
$\mathsf C_i^aF^{a,b}=F^{a,b}=\mathsf C_q^bF^{a,b}$. Hence
\[
 \Theta^{(i)}(F^{a,b})\ge0,\qquad
 \Theta^{(q)}(F^{a,b})\ge0.
\]
Discard the nonnegative $i$ term in telescoping to obtain
\begin{align}
 0\le\alpha_q\Theta^{(q)}(F^{a,b})
     &\le|\mathcal R(F^{a,b})|+
                         \alpha_p|\Theta^{(p)}(F^{a,b})|\notag\\
     &\le D+\alpha_p|\Theta^{(p)}(F^{a,b})|.
                                      \label{eq:cube_second_reduction}
\end{align}
For the third Cauchy--Schwarz step define
\[
 F^{a,b,c}=\mathsf C_p^cF^{a,b}\quad(c=0,1).
\]
Then
\begin{equation}\label{eq:cube_third_cs}
 |\Theta^{(p)}(F^{a,b})|^2
       \le\Theta^{(p)}(F^{a,b,0})\Theta^{(p)}(F^{a,b,1}).
\end{equation}
The tuple $F^{a,b,c}$ is fixed by a copy in every coordinate, so
all three forms $\Theta^{(m)}(F^{a,b,c})$ are nonnegative.
Telescoping now has no negative terms to control and gives
\[
 0\le\alpha_p\Theta^{(p)}(F^{a,b,c})
   \le\sum_{m=1}^3\alpha_m\Theta^{(m)}(F^{a,b,c})
   =\mathcal R(F^{a,b,c})\le D.
\]
Substituting this bound in \eqref{eq:cube_third_cs} and taking
nonnegative square roots yields
\[
 |\Theta^{(p)}(F^{a,b})|\le D/\alpha_p.
\]
Equation~\eqref{eq:cube_second_reduction} then gives
$0\le\Theta^{(q)}(F^{a,b})\le2D/\alpha_q$, and
\eqref{eq:cube_second_cs} gives
\[
 |\Theta^{(q)}(F^a)|\le2D/\alpha_q.
\]
This holds for each of the two $q\ne i$. Therefore
\eqref{eq:cube_first_reduction} gives
$0\le\Theta^{(i)}(F^a)\le5D/\alpha_i$, and
\eqref{eq:cube_first_cs} finally gives
\[
 |\Theta^{(i)}(F)|\le5D/\alpha_i.
\]
Restore the factor $\prod_jm_j$ and absorb $5/\alpha_i$ into $C_\alpha$.
\end{proof}

\begin{proposition}[Edge tree estimate]\label{prop:edge_tree}
If $F_{000}=F_{001}=f\in C_b(E_3;\R)$ and all other slots are $1$,
then $\Theta^{(1)}_\T(F)=\Theta^{(2)}_\T(F)=0$ and
\[
 0\le\Theta^{(3)}_{\T,\lambda,r}(F)
 \le C_\alpha\lambda^{37}R(r)^{21}|Q_\T|M_{\T,2}(f)^2.
\]
\end{proposition}
\begin{proof}
The input product is
\[
 \mathcal F_F(x)=f(x_1^0,x_2^0,x_3^0)
                         f(x_1^0,x_2^0,x_3^1).
\]
It is independent of $x_1^1,x_2^1$. By absolute integrability we
may integrate the corresponding kernel factors first, and
\[
 \int_\R H_1(x_1^1)\dd x_1^1=\int_\R\h=0,\qquad
 \int_\R H_2(x_2^1)\dd x_2^1=\int_\R\h=0.
\]
Thus $\Theta^{(1)}_\T(F)=\Theta^{(2)}_\T(F)=0$, and the
telescoping identity becomes
\[
 \alpha_3\Theta^{(3)}_\T(F)=\mathcal R_\T(F).
\]
The tuple satisfies $F=\mathsf C_3^0F=\mathsf C_3^1F$, so the
positivity conclusion of Lemma~\ref{lem:cube_cs} applies.
Integrating $G_1(x_1^1),G_2(x_2^1)$, each of mass one,
gives the nonnegative square representation
\[
 \Theta^{(3)}_\T(F)
 =\int_{\Omega_\T}\int_{E_2}G_1(x_1)G_2(x_2)
       \left(\int_\R f(x_1,x_2,v)H_3(v)\dd v\right)^2
                     \dd x_1\dd x_2\,d\nu(p,t).
\]
The edge remainder estimate now gives
\[
 0\le\Theta^{(3)}_\T(F)
    \le\frac{D_{\lambda,r}}{\alpha_3}|Q_\T|
               M_{\T,2}(f)M_{\T,2}(f).
\]
Absorb $\alpha_3^{-1}$ into $C_\alpha$ to conclude.
\end{proof}

\section{The model form and its localization}

\begin{definition}[Truncated model operator and form]\label{def:model}
Let $c:(0,\infty)\to\C$ be measurable with $|c(t)|\le1$, let
$u\in E_3$, and let $0<a<b<\infty$. The inputs in this definition
are Schwartz functions $f_j\in\Sp(E_3;\C)$. Put
\[
 \kappa_1=\kappa_2=\g*\varphi,\qquad\kappa_3=\h*\h*\psi,
\]
\begin{align*}
 U^{a,b}_{u,c}(f_1,f_2,f_3)(x)&=\int_a^bc(t)
    \prod_{j=1}^3(f_j*_j(\kappa_j)_{t^{\alpha_j},t^{\alpha_j}u_j})(x)
                  \frac{\dd t}{t},\\
 \Lambda^{a,b}_{u,c}(f)&=\int f_0(x)U^{a,b}_{u,c}(f_1,f_2,f_3)(x)\dd x.
\end{align*}
The coordinate convolution is
\begin{align*}
 &(f_j*_j(\kappa_j)_{t^{\alpha_j},t^{\alpha_j}u_j})(x)\\
 &\quad=\int_\R f_j(x_1,\ldots,x_{j-1},y_j,x_{j+1},\ldots,x_3)
    (\kappa_j)_{t^{\alpha_j},t^{\alpha_j}u_j}(x_j-y_j)\dd y_j.
\end{align*}
The translated factors can equivalently be written
\[
 (\g*\varphi_{1,u_j})_{(t^{\alpha_j})}\quad(j=1,2),
 \qquad (\h*\h*\psi_{1,u_3})_{(t^{\alpha_3})}\quad(j=3).
\]
For Schwartz inputs, $\Lambda_{u,c}$ denotes the full-scale integral,
whose convergence is proved in Lemma~\ref{lem:model_convergence}.
An almost-everywhere bound on $c$ can be made pointwise by changing
$c$ on a null set.
\end{definition}

\begin{definition}[Local form and energies]\label{def:local_model}
For a finite collection $\mathcal Q$ and bounded continuous inputs
$f_0,f_1,f_2,f_3$ define
\begin{align*}
 L_{\mathcal Q,u,c}(f)=-\int_{\Omega_{\mathcal Q}}c(t)\int_{E_6}
 &f_0(x_1,x_2,x_3)f_1(y_1,x_2,x_3)
 f_2(x_1,y_2,x_3)f_3(x_1,x_2,y_3)\\
 &\times\varphi_{t^{\alpha_1},p_1+t^{\alpha_1}u_1}(x_1)
                 \g_{t^{\alpha_1},p_1}(y_1)\\
 &\times\varphi_{t^{\alpha_2},p_2+t^{\alpha_2}u_2}(x_2)
                 \g_{t^{\alpha_2},p_2}(y_2)\\
 &\times\h_{t^{\alpha_3},p_3}(x_3)
       (\h*\psi_{1,-u_3})_{t^{\alpha_3},p_3}(y_3)
                  \dd x\dd y\,d\nu(p,t).
\end{align*}
For real inputs set, with $z=(x_1,x_2,y_1,y_2,r)$,
\begin{align*}
 A_f(p,t,z)&=\int_\R f_0(x_1,x_2,v)f_1(y_1,x_2,v)f_2(x_1,y_2,v)
                        \h_{t^{\alpha_3},p_3}(v)\dd v,\\
 B_f(p,t,z)&=\int_\R f_3(x_1,x_2,v)
                        \h_{t^{\alpha_3},p_3+rt^{\alpha_3}}(v)\dd v,
\end{align*}
and let $d\mu_{p,t,u}(z)$ be
\[
 |\varphi_{t^{\alpha_1},p_1+t^{\alpha_1}u_1}(x_1)|\g_{t^{\alpha_1},p_1}(y_1)
 |\varphi_{t^{\alpha_2},p_2+t^{\alpha_2}u_2}(x_2)|\g_{t^{\alpha_2},p_2}(y_2)
 |\psi(r+u_3)|\dd z.
\]
Define the nonnegative energies
\[
 E_{1,\mathcal Q,u}(f_0,f_1,f_2)=\int_{\Omega_{\mathcal Q}}\int A_f^2\dd\mu\,d\nu,
 \qquad
 E_{2,\mathcal Q,u}(f_3)=\int_{\Omega_{\mathcal Q}}\int B_f^2\dd\mu\,d\nu.
\]
Here $\dd z=\dd x_1\dd x_2\dd y_1\dd y_2\dd r$.
The second definition also applies to a real measurable input $f_3$
whose displayed fiber integrals exist almost everywhere, and to an
arbitrary measurable region $V\subset E_3\times(0,\infty)$:
\[
 E_{2,V,u}(f_3)=\int_V\int_{E_5}|B_f(p,t,z)|^2
                         \dd\mu_{p,t,u}(z)\,d\nu(p,t)\in[0,\infty].
\]
This is initially a nonnegative extended integral; for $L^2$ inputs
the fiber integrals exist almost everywhere, as verified in
Lemma~\ref{lem:global_energy}. Changing their values on the exceptional
null set does not change the energy.
\end{definition}

\begin{lemma}[Localization sign and integrability]\label{lem:localization}
The local forms and energies are finite on bounded continuous inputs
for finite collections. On Schwartz inputs, replacing the region in
$L$ by $E_3\times(a,b)$ gives exactly $\Lambda^{a,b}_{u,c}$.
Exhausting the center variable by bounded sets at these fixed scales
is justified by absolute integrability of the expanded integral.
\end{lemma}
\begin{proof}
For each $p,t$, translation invariance and normalized dilation give
\[
 \int_{E_5}\dd\mu_{p,t,u}
   =\|\varphi\|_1^2\|\g\|_1^2\|\psi\|_1=:C_\mu,
 \qquad
 \nu(\Omega_{\mathcal Q})=(\log2)\sum_{Q\in\mathcal Q}|Q|<\infty.
\]
The second equality follows from disjointness of the box-scale
regions. In particular
\[
 |A_f|\le\|\h\|_1\prod_{j=0}^2\|f_j\|_\infty,
 \qquad |B_f|\le\|\h\|_1\|f_3\|_\infty,
\]
and hence
\begin{align*}
 E_{1,\mathcal Q,u}(f_0,f_1,f_2)
 &\le C_\mu\|\h\|_1^2\nu(\Omega_{\mathcal Q})
                                     \prod_{j=0}^2\|f_j\|_\infty^2,\\
 E_{2,\mathcal Q,u}(f_3)
 &\le C_\mu\|\h\|_1^2\nu(\Omega_{\mathcal Q})\|f_3\|_\infty^2.
\end{align*}
The absolute integral defining $L_{\mathcal Q,u,c}$ is at most
\[
 \nu(\Omega_{\mathcal Q})\prod_{j=0}^3\|f_j\|_\infty\,
 \|\varphi\|_1^2\|\g\|_1^2\|\h\|_1\|\h*\psi\|_1.
\]
These estimates prove local finiteness. They also justify the
expansion of either square by Fubini, since replacing its two
copies of $\h$ by $|\h|$ gives the same displayed finite bounds.
All the functions under consideration are measurable; their fiber
integrals are measurable by the parameter-integral theorem.
For the global-center assertion, bound $f_1,f_2,f_3$ by their
uniform norms and retain $|f_0(x)|$. At each fixed $t>0$, Tonelli
allows the absolute kernel factors to be integrated first in
$y_1,y_2,y_3$ and then in $p_1,p_2,p_3$. Their $L^1$ norms are
invariant under the normalized dilations and translations, giving
the bound
\[
 \|f_0\|_1\prod_{j=1}^3\|f_j\|_\infty\,
 \|\varphi\|_1^2\|\g\|_1^2\|\h\|_1\|\h*\psi\|_1
\]
for the absolute spatial integral. Lemma~\ref{lem:schwartz} makes
every factor finite. Integrating in $t\in[a,b]$ multiplies this
bound by $\log(b/a)$, so the entire expanded integrand is absolutely
integrable at these fixed scales. This proves Fubini and dominated
center exhaustion directly from the Schwartz assumptions.
For a one-dimensional kernel $v$, write $\check v(w)=v(-w)$.
For Schwartz kernels $a,b$, the change of variables $q=x-p$ gives
\[
 \int_\R a(x-p)b(y-p)\dd p=(a*\check b)(x-y).
\]
In coordinates $j=1,2$, apply this with
$a=\varphi_{s,su_j}$, $b=\g_{(s)}$, and $s=t^{\alpha_j}$.
Since $\check\g=\g$, the result is
$(\g*\varphi)_{s,su_j}(x_j-y_j)$.
In the third coordinate, evenness of $\psi$ and oddness of $\h$ imply
\[
 \check{(\psi_{1,-u_3})}=\psi_{1,u_3},\qquad
 \check{(\h*\psi_{1,-u_3})}
       =-\h*\psi_{1,u_3}.
\]
The same center identity therefore gives
\[
 \int_\R\h_s(x-p)(\h*\psi_{1,-u_3})_s(y-p)\dd p
   =-(\h*\h*\psi)_{s,su_3}(x-y).
\]
The outer minus sign cancels this minus sign, giving the model in
Definition~\ref{def:model}. Also $U(u_1,u_2,-u_3)=U(u)$.
\end{proof}

\begin{lemma}[First local Cauchy--Schwarz]\label{lem:local_cs}
For real inputs and possibly complex $c$,
\[
 |L_{\mathcal Q,u,c}(f)|^2\le E_{1,\mathcal Q,u}(f_0,f_1,f_2)
                                    E_{2,\mathcal Q,u}(f_3).
\]
\end{lemma}
\begin{proof}
The convolution definition and the change of variables from the
physical translation to $r$ give
\[
 (\h*\psi_{1,-u_3})_{s,p}(y)
  =s^{-1}\int_\R\h((y-p)/s-r)\psi(r+u_3)\dd r
  =\int_\R\h_{s,p+rs}(y)\psi(r+u_3)\dd r.
\]
Define the signed density
\begin{align*}
 w_{p,t,u}(z)={}&
 \varphi_{t^{\alpha_1},p_1+t^{\alpha_1}u_1}(x_1)
            \g_{t^{\alpha_1},p_1}(y_1)\\
 &\times\varphi_{t^{\alpha_2},p_2+t^{\alpha_2}u_2}(x_2)
            \g_{t^{\alpha_2},p_2}(y_2)\psi(r+u_3).
\end{align*}
Thus $\dd\mu_{p,t,u}=|w_{p,t,u}(z)|\dd z$. Substitution of the
convolution expansion and integration first in $x_3,y_3$ give the
exact factorization
\[
 L_{\mathcal Q,u,c}(f)=-\int_{\Omega_{\mathcal Q}}c(t)
       \int_{E_5}A_f(p,t,z)B_f(p,t,z)w_{p,t,u}(z)\dd z\,d\nu(p,t).
\]
Absolute integrability follows from Lemma~\ref{lem:localization},
also after expanding the extra $r$ integral using
$\|\h\|_1\|\psi\|_1$ in place of $\|\h*\psi\|_1$.
The triangle inequality and $|c|\le1$ yield
\begin{align*}
 |L_{\mathcal Q,u,c}(f)|&\le
 \int_{\Omega_{\mathcal Q}}\int_{E_5}|A_fB_f|\dd\mu\,d\nu\\
 &\le\left(\int_{\Omega_{\mathcal Q}}\int A_f^2\dd\mu\,d\nu\right)^{1/2}
     \left(\int_{\Omega_{\mathcal Q}}\int B_f^2\dd\mu\,d\nu\right)^{1/2}.
\end{align*}
The last step is Cauchy--Schwarz on the single positive measure
space with measure $\one_{\Omega_{\mathcal Q}}\dd\mu\,d\nu$,
including both the center and the scale variables. Squaring proves
the assertion.
\end{proof}

\begin{lemma}[Local energy bounds]\label{lem:local_energies}
For finite convex trees and real bounded continuous inputs,
\begin{align}
 E_{1,\T,u}(f_0,f_1,f_2)&\le C U(u)^{50}|Q_\T|
                              \prod_{j=0}^2M_{\T,4}(f_j)^2,\label{eq:energy1}\\
 E_{2,\T,u}(f_3)&\le C U(u)^{100}|Q_\T|M_{\T,2}(f_3)^2.\label{eq:energy2}
\end{align}
\end{lemma}
\begin{proof}
Write $s_i=t^{\alpha_i}$ and set
\begin{align*}
 W_{p,t,u}(x_1,x_2,y_1,y_2)={}&
 |\varphi_{s_1,p_1+s_1u_1}(x_1)|\g_{s_1,p_1}(y_1)
 |\varphi_{s_2,p_2+s_2u_2}(x_2)|\g_{s_2,p_2}(y_2),\\
 D_{\lambda,p,t}(x_1,x_2,y_1,y_2)={}&
 \prod_{i=1}^2\g_{\lambda s_i,p_i}(x_i)
                         \g_{\lambda s_i,p_i}(y_i).
\end{align*}
The unscaled superposition satisfies
\begin{align*}
 \int_1^\infty\prod_{a=1}^4\g_{(\lambda)}(z_a)
                          \lambda^{-47}\dd\lambda
 &=\int_1^\infty e^{-\pi|z|^2/\lambda^2}\lambda^{-51}\dd\lambda\\
 &\ge e^{-\pi}\int_{1+|z|}^{2(1+|z|)}\lambda^{-51}\dd\lambda
   =\frac{e^{-\pi}(1-2^{-50})}{50}(1+|z|)^{-50}.
\end{align*}
Order $50$ Schwartz decay of the four original factors and
\[
 1+|z|\le (1+|z-(u_1,0,u_2,0)|)U(u)
\]
give their upper bound
$C U(u)^{50}(1+|z|)^{-50}$, as in Lemma~\ref{lem:domination}.
Substitute
\[
 z=\big((x_1-p_1)/s_1,(y_1-p_1)/s_1,
        (x_2-p_2)/s_2,(y_2-p_2)/s_2\big)
\]
and multiply both sides by $s_1^{-2}s_2^{-2}$. The normalized result is
\begin{equation}\label{eq:scaled_superposition}
 W_{p,t,u}\le C U(u)^{50}\int_1^\infty
                     D_{\lambda,p,t}\lambda^{-47}\dd\lambda.
\end{equation}

In $E_1$, the factor $A_f$ is independent of $r$, so first use
\[
 \int_\R|\psi(r+u_3)|\dd r=\|\psi\|_1.
\]
Apply \eqref{eq:scaled_superposition} to the nonnegative integrand
$A_f^2W_{p,t,u}$ before expanding the square. With
$\dd x'\dd y'=\dd x_1\dd x_2\dd y_1\dd y_2$, this gives
\[
 E_{1,\T,u}\le C U(u)^{50}\int_1^\infty
 \left(\int_{\Omega_\T}\int_{E_4}A_f(p,t,z)^2
                  D_{\lambda,p,t}(x',y')\dd x'\dd y'\,d\nu\right)
                        \lambda^{-47}\dd\lambda.
\]
For each fixed $\lambda$ the inner integral is finite by boundedness
of the inputs and integrability of the Gaussian factors.
Expanding the square gives
\begin{align*}
 A_f(p,t,z)^2=\int_{\R^2}
 &f_0(x_1,x_2,v)f_0(x_1,x_2,w)
  f_1(y_1,x_2,v)f_1(y_1,x_2,w)\\
 &\times f_2(x_1,y_2,v)f_2(x_1,y_2,w)
                  \h_{s_3,p_3}(v)\h_{s_3,p_3}(w)\dd v\dd w.
\end{align*}
Under the measure-preserving coordinate identification
\[
 (x_1^0,x_1^1,x_2^0,x_2^1,x_3^0,x_3^1)
                    =(x_1,y_1,x_2,y_2,v,w),
\]
the resulting tuple is
\[
 F^{[1]}_{00b}=f_0,\quad F^{[1]}_{10b}=f_1,\quad
 F^{[1]}_{01b}=f_2,\quad F^{[1]}_{11b}=1\quad(b=0,1),
\]
and the inner integral is exactly
$\Theta^{(3)}_{\T,\lambda,0}(F^{[1]})$. Consequently
\begin{equation}\label{eq:first_energy_cube_comparison}
 E_{1,\T,u}\le C U(u)^{50}\int_1^\infty
       \Theta^{(3)}_{\T,\lambda,0}(F^{[1]})\lambda^{-47}\dd\lambda.
\end{equation}
The constant slots have $M_{\T,4}(1)=(2/9)^{3/4}\le1$.
Each of $f_0,f_1,f_2$ appears twice, and hence
Proposition~\ref{prop:cube_tree} gives
\[
 E_{1,\T,u}\le C U(u)^{50}|Q_\T|
        \prod_{j=0}^2M_{\T,4}(f_j)^2
          \int_1^\infty\lambda^{37}\lambda^{-47}\dd\lambda.
\]
The remaining integral is $\int_1^\infty\lambda^{-10}\dd\lambda=1/9$,
which proves \eqref{eq:energy1}.

For $E_2$, use \eqref{eq:scaled_superposition} and
\[
 |\psi(r+u_3)|\le C(1+|u_3|)^{50}(1+|r|)^{-50}
                  \le C U(u)^{50}R(r)^{-50}.
\]
Both estimates are applied while $B_f^2$ is still a nonnegative
square. For fixed $\lambda,r$, its expansion is
\begin{align*}
 B_f(p,t,z)^2=\int_{\R^2}
 &f_3(x_1,x_2,v)f_3(x_1,x_2,w)\\
 &\times\h_{s_3,p_3+rs_3}(v)\h_{s_3,p_3+rs_3}(w)\dd v\dd w.
\end{align*}
With the same coordinate identification as above, multiplication by
$D_{\lambda,p,t}$ and integration on $\Omega_\T\times E_4$ give
exactly $\Theta^{(3)}_{\T,\lambda,r}(F^{[2]})$, where
\[
 F^{[2]}_{000}=F^{[2]}_{001}=f_3,\qquad
 F^{[2]}_j=1\quad\text{for }j\in J\setminus\{000,001\}.
\]
Thus
\begin{equation}\label{eq:second_energy_cube_comparison}
 E_{2,\T,u}\le C U(u)^{100}\int_1^\infty\int_\R
       \Theta^{(3)}_{\T,\lambda,r}(F^{[2]})
                      \lambda^{-47}R(r)^{-50}\dd r\dd\lambda.
\end{equation}
Proposition~\ref{prop:edge_tree} now yields
\begin{align*}
 E_{2,\T,u}\le{}& C U(u)^{100}|Q_\T|M_{\T,2}(f_3)^2\\
 &\times\int_1^\infty\lambda^{37}\lambda^{-47}\dd\lambda
        \int_\R R(r)^{21}R(r)^{-50}\dd r.
\end{align*}
The two integrals are $1/9$ and
$\int_\R(1+|r|)^{-29}\dd r=1/14$, respectively, proving
\eqref{eq:energy2}. All superpositions above involve nonnegative
square integrands, so Tonelli applies even before these bounds prove
finiteness of the $\lambda,r$ integrals. Expansion of a square is
made only at fixed $\lambda,r$, where the absolute integrability
argument from Lemma~\ref{lem:localization} applies.
\end{proof}

\begin{lemma}[Global second energy]\label{lem:global_energy}
For real $f\in L^2(E_3)$,
\begin{equation}\label{eq:global_energy}
 E_{2,E_3\times(0,\infty),u}(f)
 =\|\varphi\|_1^2\|\psi\|_1\frac{\pi}{\alpha_3}\|f\|_2^2.
\end{equation}
In particular the same upper bound holds on every measurable region,
and the energies add over disjoint regions.
For a finite convex tree and real $f\in C_b(E_3)\cap L^2(E_3)$,
combining this with the local bound gives
\begin{equation}\label{eq:second_energy_minimum}
 E_{2,\T,u}(f)\le C\min\big\{
        U(u)^{100}|Q_\T|M_{\T,2}(f)^2,\ \|f\|_2^2\big\}.
\end{equation}
\end{lemma}
\begin{proof}
Define the fiber transform, including its center variable, by
\begin{equation}\label{eq:fiber_gaussian_transform}
 H_t f(x_1,x_2,q)=\int_\R f(x_1,x_2,v)
                              \h_{t^{\alpha_3},q}(v)\dd v.
\end{equation}
For Schwartz $f$ this is an ordinary absolutely integrable fiber
integral. For $f\in L^2(E_3)$, Fubini gives
$f(x_1,x_2,\cdot)\in L^2(\R)$ for almost every $(x_1,x_2)$.
On such a fiber, Cauchy--Schwarz makes \eqref{eq:fiber_gaussian_transform}
absolutely integrable for every $q,t$, since
\[
 \|\h_{t^{\alpha_3},q}\|_2=t^{-\alpha_3/2}\|\h\|_2.
\]
Set the transform to zero on the exceptional set of fibers.

All integrands defining the energy are nonnegative. Tonelli first
integrates $y_1,y_2$ using
\[
 \int_\R\g_{t^{\alpha_i},p_i}(y_i)\dd y_i=1\quad(i=1,2),
\]
and then $p_1,p_2$ using
\[
 \int_\R|\varphi_{t^{\alpha_i},p_i+t^{\alpha_i}u_i}(x_i)|\dd p_i
                        =\|\varphi\|_1\quad(i=1,2).
\]
The remaining integrand contains
$|H_t f(x_1,x_2,p_3+rt^{\alpha_3})|^2$. For fixed $r,t$, substitute
$q=p_3+rt^{\alpha_3}$, whose Jacobian is one, and then use
$\int_\R|\psi(r+u_3)|\dd r=\|\psi\|_1$. We obtain the exact identity
\begin{equation}\label{eq:global_energy_fiber_reduction}
 E_{2,E_3\times(0,\infty),u}(f)
  =\|\varphi\|_1^2\|\psi\|_1
    \int_{E_2}\int_0^\infty\int_\R
       |H_t f(x_1,x_2,q)|^2\dd q\,\frac{\dd t}{t}\dd x_1\dd x_2.
\end{equation}
This equality holds for nonnegative extended integrals before any
finiteness is known.

Fix a fiber $F(v)=f(x_1,x_2,v)\in L^2(\R)$ and put $s=t^{\alpha_3}$.
Oddness of $\h$ gives
\[
 H_t f(x_1,x_2,q)=-(F*\h_{(s)})(q),\qquad
 \widehat{H_t f(x_1,x_2,\cdot)}(\xi)
                 =-\widehat F(\xi)\widehat\h(s\xi).
\]
The Fourier identity is first the Schwartz convolution identity and
then its $L^2$ extension; convolution with $\h_{(s)}\in L^1$
is bounded on $L^2$. Plancherel and Tonelli give
\begin{align*}
 \int_0^\infty\int_\R|H_t f(x_1,x_2,q)|^2\dd q\,\frac{\dd t}{t}
 &=\int_0^\infty\int_\R|\widehat F(\xi)|^2
              |\widehat\h(t^{\alpha_3}\xi)|^2\dd\xi\,\frac{\dd t}{t}\\
 &=\int_\R|\widehat F(\xi)|^2
     \left(\int_0^\infty|\widehat\h(t^{\alpha_3}\xi)|^2
                                          \frac{\dd t}{t}\right)\dd\xi.
\end{align*}
For $\xi\ne0$, put $v=t^{\alpha_3}|\xi|$, so that
$\dd t/t=\dd v/(\alpha_3v)$. Since
$|\widehat\h(\eta)|^2=4\pi^2\eta^2e^{-2\pi\eta^2}$, the multiplier
integral is explicitly
\[
 \int_0^\infty|\widehat\h(t^{\alpha_3}\xi)|^2\frac{\dd t}{t}
 =\frac{4\pi^2}{\alpha_3}\int_0^\infty v e^{-2\pi v^2}\dd v
 =\frac{\pi}{\alpha_3}\quad(\xi\ne0).
\]
At $\xi=0$ that integral is zero; this single frequency is a null set.
Plancherel once more therefore gives, on almost every fiber,
\[
 \int_0^\infty\int_\R|H_t f(x_1,x_2,q)|^2\dd q\,\frac{\dd t}{t}
 =\frac{\pi}{\alpha_3}\int_\R|\widehat F(\xi)|^2\dd\xi
 =\frac{\pi}{\alpha_3}\int_\R|f(x_1,x_2,v)|^2\dd v.
\]
Insert this in \eqref{eq:global_energy_fiber_reduction} and integrate
in $(x_1,x_2)$ to obtain \eqref{eq:global_energy}.
Restriction to a measurable region decreases the nonnegative
integral, and disjoint regions have additive indicator functions,
which proves both remaining assertions.
Finally, apply \eqref{eq:energy2} as well as the global bound and
enlarge the constant to the maximum of their two constants. This
proves \eqref{eq:second_energy_minimum}.
\end{proof}

\begin{corollary}[Local model estimate]\label{cor:local_model}
\[
 |L_{\T,u,c}(f)|\le C U(u)^{75}|Q_\T|
          \left(\prod_{j=0}^2M_{\T,4}(f_j)\right)M_{\T,2}(f_3).
\]
\end{corollary}
\begin{proof}
Take the square root of the product of \eqref{eq:energy1} and
\eqref{eq:energy2}. The translation exponent is $(50+100)/2=75$.
\end{proof}

\section{Stopping time and the initial exponent range}

\begin{exttheorem}[One-dimensional maximal theorem and differentiation]
\label{ext:maximal}
For a locally integrable function on $\R$, define
\[
 Mh(x)=\sup_{r>0}\frac1{2r}\int_{x-r}^{x+r}|h(y)|\dd y.
\]
For every $1<q<\infty$ there is $C_q$ such that
$\|Mh\|_q\le C_q\|h\|_q$. Lebesgue differentiation holds for
locally integrable functions, also along the nested standard dyadic
intervals containing a point. The corresponding fiber statements hold
for functions on $E_3$ whose fibers satisfy the respective hypotheses
almost everywhere.
\end{exttheorem}
The fiber statements follow by Fubini. Fourier inversion, Plancherel,
H\"older, dominated convergence, and elementary measure theory are
background results.

\begin{lemma}[Maximal sizes]\label{lem:maximal_size}
For $1\le s<q<\infty$ and $g\in L^q(E_3)$ define
\[
 \mathfrak M_sg(x)=\sup_{t>0}(|g|^s*\Theta_t(x))^{1/s}.
\]
This is measurable, $\|\mathfrak M_sg\|_q\le C_{s,q}\|g\|_q$,
and, for every box $Q$ and $x\in Q$,
\begin{equation}\label{eq:ancestor_size}
 \mathfrak M_sg(x)\ge(2/3)^{30/s}M_{\{Q\},s}(g).
\end{equation}
\end{lemma}
\begin{proof}
Annular decomposition of $s^{-1}(1+|x|/s)^{-10}$ gives domination
of convolution by a fixed multiple of the one-dimensional maximal
operator. Applying this in the three coordinates yields
\[
 |g|^s*\Theta_t\le C M^{(1)}M^{(2)}M^{(3)}(|g|^s).
\]
Apply External Theorem~\ref{ext:maximal} three times in $L^{q/s}$.
The convolutions depend continuously on $(x,t)$ for $t>0$, by
H\"older and continuity of translations and dilations of the kernel
in $L^{(q/s)'}$. Thus the supremum may be taken over rational $t>0$.
Finally $|x_i-c(Q)_i|\le\ell(Q)^{\alpha_i}/2$ implies
\[
 \Theta_{\ell(Q)}(x-y)\ge(2/3)^{30}
                      \Theta_{\ell(Q)}(c(Q)-y).
\]
Integrate this inequality to get \eqref{eq:ancestor_size}.
\end{proof}

\begin{definition}[Stopping data]\label{def:stopping}
Let $d\ge2$, let $\mathcal Q_0$ be a finite convex collection, and let
$g_j\in L^{q_j}(E_3)$ be nonzero as $L^{q_j}$ classes, where
$1\le s_j<q_j<\infty$ for $0\le j<d$.
For $Q\in\mathcal Q_0$ set
\begin{equation}\label{eq:stopping_data}
 a_j(Q)=\max_{\substack{Q'\in\mathcal Q_0\\Q\subseteq Q'}}
                         M_{\{Q'\},s_j}(g_j).
\end{equation}
For $n=(n_j)_{0\le j<d}\in\Z^d$, define
\begin{align}
 \mathcal P_n
 &=\{Q\in\mathcal Q_0:
       \forall j<d,\ 2^{n_j-1}<a_j(Q)\le2^{n_j}\},
                  \label{eq:stopping_levels}\\
 \mathcal P_n^{\max}
 &=\{Q\in\mathcal P_n:
       \nexists R\in\mathcal P_n\text{ with }Q\subsetneq R\},
                  \label{eq:stopping_roots}\\
 \T_Q&=\{Q'\in\mathcal P_n:Q'\subseteq Q\}
              \qquad(Q\in\mathcal P_n^{\max}).
                  \label{eq:stopping_trees}
\end{align}
The forest is the family indexed by pairs $(n,Q)$ with
$Q\in\mathcal P_n^{\max}$; its root is $Q$.
\end{definition}

\begin{lemma}[Stopping-time sum]\label{lem:stopping}
Suppose also that $\sum_{j<d}q_j^{-1}=1$. The $\T_Q$ partition
$\mathcal Q_0$ into convex trees, and
\begin{equation}\label{eq:stopping_sum}
 \sum_n\sum_{Q\in\mathcal P_n^{\max}}|Q|
             \prod_{j<d}M_{\T_Q,s_j}(g_j)
 \le C_{s,q,d}\prod_{j<d}\|g_j\|_{q_j}.
\end{equation}
Only finitely many summands are nonzero for a fixed $\mathcal Q_0$.
\end{lemma}
\begin{proof}
The ancestor set in \eqref{eq:stopping_data} is nonempty because
it contains $Q$. Each of its sizes is finite by H\"older, since
$|g_j|^{s_j}\in L^{q_j/s_j}$ and the fixed translate of
$\Theta_{\ell(Q')}$ belongs to the conjugate space. It is strictly
positive because $g_j$ is nonzero almost everywhere on a set of
positive measure and $\Theta$ is strictly positive. A maximum over
the finite nonempty ancestor set is therefore finite and positive.
For every positive real number $a$ there is exactly one $n\in\Z$
with $2^{n-1}<a\le2^n$. Applying this separately to every $a_j(Q)$
places each $Q$ in exactly one $\mathcal P_n$.

If $Q\subseteq R$, then
\[
 \{R'\in\mathcal Q_0:R\subseteq R'\}
 \subseteq\{Q'\in\mathcal Q_0:Q\subseteq Q'\},
 \qquad a_j(Q)\ge a_j(R).
\]
For $Q_1\subseteq Q_2\subseteq Q_3$ in $\mathcal Q_0$ with
$Q_1,Q_3\in\mathcal P_n$, one consequently has
\[
 2^{n_j-1}<a_j(Q_3)\le a_j(Q_2)\le a_j(Q_1)\le2^{n_j}.
\]
Thus $Q_2\in\mathcal P_n$, proving convexity along dyadic chains.
Intersecting anisotropic dyadic boxes are nested. Distinct maximal
members of $\mathcal P_n$ are therefore disjoint, and finiteness
gives a unique maximal ancestor for every member. The trees
$\T_Q$ are convex, partition $\mathcal P_n$, and, as $n$ varies,
partition $\mathcal Q_0$. Moreover, for $R\in\T_Q$,
$M_{\{R\},s_j}(g_j)\le a_j(R)\le2^{n_j}$, whence
\begin{equation}\label{eq:stopping_tree_size}
 M_{\T_Q,s_j}(g_j)\le2^{n_j}.
\end{equation}

Set
\begin{equation}\label{eq:stopping_level_constant}
 c_*=\frac12\min_{0\le j<d}(2/3)^{30/s_j}>0,
 \qquad E_{j,n}=\{x:\mathfrak M_{s_j}g_j(x)\ge c_*2^n\}.
\end{equation}
The sets $E_{j,n}$ are measurable by Lemma~\ref{lem:maximal_size}.
Fix a root $Q\in\mathcal P_n^{\max}$ and $j<d$. An ancestor
$R\supseteq Q$ realizes the maximum in $a_j(Q)$, so
$M_{\{R\},s_j}(g_j)>2^{n_j-1}$. For $x\in Q\subseteq R$,
the coordinate inequalities
\[
 |x_i-c(R)_i|\le\tfrac12\ell(R)^{\alpha_i},\qquad
 1+\ell(R)^{-\alpha_i}|x_i-y_i|
 \le\tfrac32\big(1+\ell(R)^{-\alpha_i}|c(R)_i-y_i|\big)
\]
give
\[
 \Theta_{\ell(R)}(x-y)\ge(2/3)^{30}
                              \Theta_{\ell(R)}(c(R)-y).
\]
After integration and taking the $s_j$th root,
\[
 \mathfrak M_{s_j}g_j(x)
 \ge(2/3)^{30/s_j}M_{\{R\},s_j}(g_j)>c_*2^{n_j}.
\]
Thus $Q\subseteq E_{j,n_j}$ for every $j$, and disjointness of
roots at the same level vector gives
\begin{equation}\label{eq:stopping_root_measure}
 \sum_{Q\in\mathcal P_n^{\max}}|Q|\le |E_{j,n_j}|.
\end{equation}

For $j<d$, define $L_j=\|g_j\|_{q_j}>0$ and
$\sigma_j=\log_2L_j$, so $2^{\sigma_j}=L_j$. Define the
tie-breaking selector and its level sets:
\begin{align*}
 J(n)&=\min\left\{j<d:q_j(n_j-\sigma_j)
                    =\max_{k<d}q_k(n_k-\sigma_k)\right\},\\
 \mathcal N_j&=\{n\in\Z^d:J(n)=j\},\\
 S_j&=\sum_{n\in\mathcal N_j}2^{\sum_{k<d}n_k}
                       \sum_{Q\in\mathcal P_n^{\max}}|Q|.
\end{align*}
The $\mathcal N_j$ partition $\Z^d$. The finite maximum exists,
so $J$ is defined for every vector. If $n\in\mathcal N_j$ and
$k\ne j$, then
\begin{equation}\label{eq:stopping_dominant}
 q_j(n_j-\sigma_j)\ge q_k(n_k-\sigma_k),\qquad
 n_k\le\sigma_k+\frac{q_j}{q_k}(n_j-\sigma_j).
\end{equation}
By \eqref{eq:stopping_root_measure}, positivity, and enlargement
of the remaining index sets,
\begin{align*}
 S_j
 &\le\sum_{n\in\mathcal N_j}2^{\sum_kn_k}|E_{j,n_j}|\\
 &\le\sum_{n_j\in\Z}2^{n_j}|E_{j,n_j}|
       \prod_{k\ne j}
       \sum_{\substack{n_k\in\Z\\
          n_k\le\sigma_k+(q_j/q_k)(n_j-\sigma_j)}}2^{n_k}\\
 &\le2^{d-1}\sum_{n_j\in\Z}2^{n_j}|E_{j,n_j}|
                 \prod_{k\ne j}
                 \left[L_k2^{(q_j/q_k)(n_j-\sigma_j)}\right].
\end{align*}
The last step is the elementary geometric-series identity
$\sum_{n\le A}2^n=2^{\lfloor A\rfloor+1}\le2^{A+1}$.
Because $\sum_{k\ne j}q_k^{-1}=1-q_j^{-1}$,
\[
 2^{n_j}\prod_{k\ne j}2^{(q_j/q_k)(n_j-\sigma_j)}
 =2^{q_jn_j-(q_j-1)\sigma_j}
 =2^{q_jn_j}L_j^{1-q_j}.
\]
Consequently
\begin{equation}\label{eq:stopping_distribution_sum}
 S_j\le2^{d-1}\left(\prod_{k\ne j}L_k\right)L_j^{1-q_j}
                 \sum_{n_j\in\Z}2^{q_jn_j}|E_{j,n_j}|.
\end{equation}
For measurable $h\ge0$ and $q>0$, Tonelli gives the discrete
layer-cake bound
\begin{align*}
 \sum_{n\in\Z}2^{qn}|\{h\ge c_*2^n\}|
 &=\int\sum_{\substack{n\in\Z\\c_*2^n\le h(x)}}2^{qn}\dd x\\
 &\le\frac{2^q}{(2^q-1)c_*^q}\int h(x)^q\dd x.
\end{align*}
For $h(x)>0$ this is a geometric sum up to
$\lfloor\log_2(h(x)/c_*)\rfloor$; for $h(x)=0$ the sum is empty.
Apply it to $h=\mathfrak M_{s_j}g_j$ and $q=q_j$, and use
Lemma~\ref{lem:maximal_size}, to obtain
\[
 \sum_{n_j\in\Z}2^{q_jn_j}|E_{j,n_j}|
 \le C\|\mathfrak M_{s_j}g_j\|_{q_j}^{q_j}
 \le C L_j^{q_j}.
\]
Equation~\eqref{eq:stopping_distribution_sum} now gives
$S_j\le C\prod_kL_k$. Finally,
\[
 \sum_n\sum_{Q\in\mathcal P_n^{\max}}|Q|
             \prod_{j<d}M_{\T_Q,s_j}(g_j)
 \le\sum_n2^{\sum_jn_j}\sum_{Q\in\mathcal P_n^{\max}}|Q|
 =\sum_{j<d}S_j\le C\prod_{j<d}\|g_j\|_{q_j}.
\]
All rearrangements above concern nonnegative series, so Tonelli
applies before finiteness of the resulting bounds is known.
\end{proof}

\begin{lemma}[Finite-forest estimates, including the endpoint]
\label{lem:forest_bound}
If $p_0,p_1,p_2>4$, $2\le p_3<\infty$, all exponents are finite,
and $\sum_jp_j^{-1}=1$, then for real Schwartz inputs
\[
 |L_{\mathcal Q_0,u,c}(f)|\le C_{\alpha,p}U(u)^{100}
                              \prod_{j=0}^3\|f_j\|_{p_j},
\]
uniformly over finite convex collections. One may replace $100$ by
$75$ when $p_3>2$, and by $25$ when $p_3=2$.
\end{lemma}
\begin{proof}
If any function is zero, the conclusion follows by multilinearity.
For $p_3>2$ apply the stopping lemma with $d=4$,
$(s_0,s_1,s_2,s_3)=(4,4,4,2)$, $q_j=p_j$, $g_j=f_j$.
Write its forest as $\mathfrak F=\{\T_Q:Q\in\mathcal P_n^{\max},
n\in\Z^4\}$. The partition of boxes gives, up to the finitely
many boundary hyperplanes of measure zero,
\[
 \Omega_{\mathcal Q_0}=\mathop{\dot\bigcup}_{\T\in\mathfrak F}
               \Omega_\T,\qquad
 L_{\mathcal Q_0,u,c}(f)=\sum_{\T\in\mathfrak F}L_{\T,u,c}(f).
\]
For each tree the Cauchy--Schwarz and energy bounds give
\begin{align*}
 |L_{\T_Q,u,c}(f)|^2
 &\le E_{1,\T_Q,u}(f_0,f_1,f_2)E_{2,\T_Q,u}(f_3),\\
 |L_{\T_Q,u,c}(f)|
 &\le C U(u)^{75}|Q|
               \left(\prod_{j=0}^2M_{\T_Q,4}(f_j)\right)
               M_{\T_Q,2}(f_3).
\end{align*}
Thus the triangle inequality and \eqref{eq:stopping_sum} imply
\begin{align*}
 |L_{\mathcal Q_0,u,c}(f)|
 &\le C U(u)^{75}\sum_n\sum_{Q\in\mathcal P_n^{\max}}
       |Q|\left(\prod_{j=0}^2M_{\T_Q,4}(f_j)\right)
                    M_{\T_Q,2}(f_3)\\
 &\le C U(u)^{75}\prod_{j=0}^3\|f_j\|_{p_j}.
\end{align*}

For $p_3=2$, choose a different stopping partition: take $d=3$,
$g_j=f_j^2$, $s_j=2$, and $q_j=p_j/2$ for $j=0,1,2$.
These satisfy $q_j>2$ and $\sum_{j=0}^2q_j^{-1}=1$.
Cauchy--Schwarz over the trees and Lemma~\ref{lem:local_cs} give
\[
 |L_{\mathcal Q_0,u,c}(f)|\le
       \left(\sum_\T E_{1,\T,u}\right)^{1/2}
       \left(\sum_\T E_{2,\T,u}\right)^{1/2}.
\]
For the first sum use \eqref{eq:energy1}, the identity
$M_{\T,4}(f_j)^2=M_{\T,2}(g_j)$, and \eqref{eq:stopping_sum}:
\begin{align*}
 \sum_{\T\in\mathfrak F}E_{1,\T,u}(f_0,f_1,f_2)
 &\le C U(u)^{50}\sum_n\sum_{Q\in\mathcal P_n^{\max}}
                |Q|\prod_{j=0}^2M_{\T_Q,4}(f_j)^2\\
 &=C U(u)^{50}\sum_n\sum_{Q\in\mathcal P_n^{\max}}
                |Q|\prod_{j=0}^2M_{\T_Q,2}(g_j)\\
 &\le C U(u)^{50}\prod_{j=0}^2\|g_j\|_{p_j/2}
  = C U(u)^{50}\prod_{j=0}^2\|f_j\|_{p_j}^2.
\end{align*}
Here the identity for the local sizes follows directly by squaring
the fourth root of the defining convolution; also
$\|f_j^2\|_{p_j/2}=\|f_j\|_{p_j}^2$. For the second sum,
nonnegativity and region disjointness give
\[
 \sum_{\T\in\mathfrak F}E_{2,\T,u}(f_3)
 =E_{2,\mathcal Q_0,u}(f_3)
 \le C\|f_3\|_2^2
\]
by the global bound \eqref{eq:global_energy}.
Taking square roots proves the endpoint with power $25$.
\end{proof}

\begin{lemma}[Truncation, exhaustion, and convergence]\label{lem:model_convergence}
For Schwartz tuples, the spatially integrated model integrand is
absolutely integrable in $t$ on $(0,\infty)$. Moreover the forest
estimate implies the identical bound for every $\Lambda^{a,b}_{u,c}$
and for $\Lambda_{u,c}$, with the same constant.
\end{lemma}
\begin{proof}
For small $t$, $\int\kappa_3=0$ and the mean value theorem give
\[
 \|f_3*_3(\kappa_3)_{t^{\alpha_3},t^{\alpha_3}u_3}\|_\infty
 \le t^{\alpha_3}\|\partial_3f_3\|_\infty
             \int|z+u_3||\kappa_3(z)|\dd z.
\]
Bound the other two convolutions by their $L^\infty$ norms and
$f_0$ in $L^1$. For $t\ge1$, the third convolution is bounded by
\[
 t^{-\alpha_3}\|\kappa_3\|_\infty
       \sup_{x_1,x_2}\int_\R|f_3(x_1,x_2,y)|\dd y.
\]
Again the other factors are bounded, and the displayed fiber supremum
is finite for Schwartz $f_3$. Thus the absolute spatial integral is
bounded by a tuple-dependent constant times
$\min(t^{\alpha_3},t^{-\alpha_3})$, integrable against $\dd t/t$.

To transfer the forest bound, replace $c$ by $c\one_{(a,b)}$.
For $N\in\N$ set
\begin{align*}
 \mathcal R_N&=\{Q(N,v):v\in\{-1,0\}^3\},\\
 B_N&=\bigcup_{R\in\mathcal R_N}R
          =\prod_{i=1}^3[-2^{N\alpha_i},2^{N\alpha_i}),\\
 \mathcal Q_N&=\{Q(k,n):-N\le k\le N,\ n\in\Z^3,
                                      \ Q(k,n)\subseteq B_N\}.
\end{align*}
At every fixed level $k$ in this range, the included boxes partition
$B_N$. Their number is finite, because each coordinate index is
bounded by the side lengths of $B_N$. Dyadic nesting proves that
each $\mathcal Q_N$ is convex, and the associated region is
\[
 \Omega_{\mathcal Q_N}=B_N\times(2^{-N-1},2^N].
\]
Indeed the level-$k$ boxes contribute
$B_N\times(2^{k-1},2^k]$. Moreover,
\[
 B_N\subseteq B_{N+1},\quad
 \mathcal Q_N\subseteq\mathcal Q_{N+1},\quad
 \Omega_{\mathcal Q_N}\subseteq\Omega_{\mathcal Q_{N+1}},\quad
 \bigcup_{N\in\N}\Omega_{\mathcal Q_N}=E_3\times(0,\infty).
\]
The union follows by choosing $N$ so that
$|p_i|<2^{N\alpha_i}$ for all $i$ and $2^{-N-1}<t\le2^N$.
At the fixed scales $[a,b]$ use the
absolute majorant of Lemma~\ref{lem:localization} to take the limit
in $L_{\mathcal Q_N,u,c\one_{(a,b)}}$. It is $\Lambda^{a,b}_{u,c}$.
Only after this center exhaustion let $a\downarrow0$, $b\uparrow\infty$,
using the scale estimate in the first paragraph. No absolute
integrability of the fully expanded six-variable kernel over all
scales is asserted or required.
\end{proof}

\begin{theorem}[Initial model range]\label{thm:initial_model}
For finite exponents $p_0,p_1,p_2>4$, $p_3\ge2$ with
$\sum_jp_j^{-1}=1$, all complex Schwartz tuples satisfy
\[
 |\Lambda^{a,b}_{u,c}(f)|+|\Lambda_{u,c}(f)|
 \le C_{\alpha,p}U(u)^{100}\prod_{j=0}^3\|f_j\|_{p_j},
\]
uniformly in $0<a<b$. Each form extends uniquely and continuously
to the product of these $L^{p_j}$ spaces.
\end{theorem}
\begin{proof}
The two preceding lemmas prove the real case. Expand each complex
input into real and imaginary parts and use the triangle inequality;
there are $2^4$ terms and each part has norm at most the original
norm. Multilinearity, the bound, and Schwartz density give the
extension by telescoping differences of tuples. The constant can be
enlarged to cover the sum of the two absolute values.
\end{proof}

\section{Fiberwise Calder\'on--Zygmund decomposition}

All operators in this section are first the fixed truncated operator
$U=U^{a,b}_{u,c}$. Every estimate is uniform in $a,b,c$.
For input exponents with $\sum_jp_j^{-1}=R^{-1}\le1$, this operator
is defined on the corresponding $L^p$ spaces by its actual integral:
fiberwise Young gives
\[
 \|f_j*_j(\kappa_j)_{t^{\alpha_j},t^{\alpha_j}u_j}\|_{p_j}
 \le\|\kappa_j\|_1\|f_j\|_{p_j}.
\]
H\"older, followed by Minkowski for $R\ge1$, then gives
\begin{align*}
 \|U(f)\|_R
 &\le\int_a^b
    \left\|\prod_{j=1}^3
      (f_j*_j(\kappa_j)_{t^{\alpha_j},t^{\alpha_j}u_j})\right\|_R
                                                     \frac{\dd t}{t}\\
 &\le\log(b/a)\prod_{j=1}^3\|\kappa_j\|_1\|f_j\|_{p_j}.
\end{align*}
This bound justifies density and good--bad identities.

\begin{definition}[Coordinate maximal functions]\label{def:fiber_maximal}
For $j\in\{1,2,3\}$, define the one-dimensional maximal operator
on the $j$th coordinate fiber by
\[
 M^{(j)}f(x)=\sup_{\rho>0}\frac1{2\rho}
  \int_{x_j-\rho}^{x_j+\rho}
  |f(x_1,\ldots,x_{j-1},y,x_{j+1},\ldots,x_3)|\dd y.
\]
On exceptional fibers where local integrability fails, assign the
value zero. For $f\in L^p(E_3)$, $1\le p<\infty$, Fubini makes
that set of fibers null; for $p>1$, External
Theorem~\ref{ext:maximal} gives
$\|M^{(j)}f\|_p\le C_p\|f\|_p$. The supremum can be restricted
to rational $\rho>0$, making the chosen representative measurable.
\end{definition}

\begin{lemma}[Kernel bounds in one fiber]\label{lem:fiber_kernel}
For $j=1,2,3$, $s>0$,
\begin{align*}
 |(\kappa_j)_{s,su_j}(y)|&\le C U(u)^{10}s^{-1}(1+|y|/s)^{-10},\\
 |\partial_y(\kappa_j)_{s,su_j}(y)|&\le C U(u)^{10}s^{-2}(1+|y|/s)^{-10},\\
 |f*_j(\kappa_j)_{s,su_j}(x)|&\le C U(u)^{10}M^{(j)}f(x).
\end{align*}
\end{lemma}
\begin{proof}
Use the order $10$ seminorms of $\kappa_j,\kappa_j'$ and
$1+|y|/s\le(1+|u_j|)(1+|y/s-u_j|)$.
For the last estimate decompose $|y|$ into the intervals
$|y|\le s$ and $2^{l-1}s<|y|\le2^ls$; the resulting coefficients
are bounded by a constant times $2^{-9l}$. Explicitly,
\begin{align*}
 |f*_j(\kappa_j)_{s,su_j}(x)|
 &\le C U(u)^{10}\left[
 s^{-1}\int_{|y|\le s}|f(x-ye_j)|\dd y\right.\\
 &\hspace{38mm}\left.
 +\sum_{l=1}^{\infty}s^{-1}2^{-10(l-1)}
                    \int_{|y|\le2^ls}|f(x-ye_j)|\dd y\right]\\
 &\le C U(u)^{10}\left(1+\sum_{l=1}^\infty2^{-9l}\right)
                                                M^{(j)}f(x).
\end{align*}
\end{proof}

\begin{lemma}[Measurable fiber decomposition]\label{lem:fiber_cz}
Let $1\le p<\infty$, $f\in L^p(E_3)$, $m\in\{1,2,3\}$, and $H>0$.
Write $x=(y,z)$ with $y=x_m$
and $z=x_{\widehat m}\in E_2$, and set
\[
 Z_f=\left\{z:\int_\R|f(y,z)|^p\dd y<\infty\right\},\qquad
 \mathcal D_1=\{[2^k n,2^k(n+1)):k,n\in\Z\}.
\]
For $I\in\mathcal D_1$, define the measurable fiber averages and
selection set by
\begin{align*}
 A_I(z)&=|I|^{-1}\int_I|f(y,z)|^p\dd y,\qquad
 f_I(z)=|I|^{-1}\int_I f(y,z)\dd y\quad(z\in Z_f),\\
 S_I&=Z_f\cap\{z:A_I(z)>H\}
       \cap\bigcap_{\substack{J\in\mathcal D_1\\I\subsetneq J}}
                                           \{z:A_J(z)\le H\},\\
 \mathcal I_z&=\{I\in\mathcal D_1:z\in S_I\}.
\end{align*}
Set $f_I=0$ off $Z_f$. On every $z\in Z_f$, the selected intervals
are disjoint and maximal among those with average exceeding $H$;
they satisfy
\[
 H<|I|^{-1}\int_I|f(y,z)|^p\dd y\le2H.
\]
Define the exceptional set, good function, and indexed bad functions by
\begin{align*}
 B&=\bigcup_{I\in\mathcal D_1}(I\times S_I),\\
 g(y,z)&=\begin{cases}
 f(y,z),&(y,z)\notin B,\ z\in Z_f,\\
 f_I(z),&y\in I,\ z\in S_I,\\
 0,&z\notin Z_f,
 \end{cases}\\
 b_I(y,z)&=\one_{S_I}(z)\one_I(y)(f(y,z)-f_I(z)),\\
 b(y,z)&=\sum_{I\in\mathcal D_1}b_I(y,z).
\end{align*}
All sets and functions in these formulas are measurable, and
$f=g+b$ almost everywhere. Then
\begin{gather*}
 \int_Ib_I(y,z)\dd y=0,\quad
 \|g\|_p\le\|f\|_p,\quad |g|\le(2H)^{1/p}\text{ a.e.},\\
 \sum_{I\in\mathcal I_z}|I|\le H^{-1}\int|f(y,z)|^p\dd y,
 \qquad \|b_I(\cdot,z)\|_1\le2(2H)^{1/p}|I|.
\end{gather*}
For finite $P\ge p$, $\|g\|_P\le(2H)^{1/p-1/P}\|f\|_p^{p/P}$.
\end{lemma}
\begin{proof}
Enumerate the countable family of dyadic intervals. Declare $I$
selected at $z$ when its $|f|^p$ average exceeds $H$ and every
proper dyadic ancestor has average at most $H$. Each condition is
measurable in $z$, so each $S_I$ and $B$ is measurable. Tonelli
shows that $Z_f$ is conull. On $Z_f$, if $J$ runs through growing
dyadic ancestors, then
\[
 A_J(z)\le |J|^{-1}\int_\R|f(y,z)|^p\dd y\longrightarrow0.
\]
Thus every interval with average above $H$ lies in a maximal one.
Two selected intervals meeting at a point are nested and hence equal.
If $I^+$ is the dyadic parent of a selected interval, then
\[
 H<A_I(z)\le2A_{I^+}(z)\le2H.
\]
Outside $\bigcup\mathcal I_z$, all dyadic interval averages
containing $y$ are at most $H$. Dyadic differentiation gives
$|f(y,z)|^p\le H$ for almost every such $y$. Fubini removes the
resulting null subset of $E_3$. On a selected interval, Jensen gives
\[
 |f_I(z)|^p\le A_I(z)\le2H,\qquad
 |I|\,|f_I(z)|^p\le\int_I|f(y,z)|^p\dd y.
\]
Summing the second inequality and retaining the unchanged complement
proves $\|g\|_p^p\le\|f\|_p^p$. Also $|g|\le(2H)^{1/p}$
almost everywhere. The defining average gives cancellation, and
H\"older gives the bad mass bound. For $z\in S_I$ the calculation is
\begin{align*}
 \int_I b_I(y,z)\dd y&=\int_I f(y,z)\dd y-|I|f_I(z)=0,\\
 \|b_I(\cdot,z)\|_1
 &\le\int_I|f(y,z)|\dd y+|I|\,|f_I(z)|\\
 &\le2|I|^{1-1/p}\left(\int_I|f(y,z)|^p\dd y\right)^{1/p}
 \le2(2H)^{1/p}|I|.
\end{align*}
For $z\notin S_I$, the function $b_I(\cdot,z)$ is zero, so both
conclusions hold there as well.
Disjointness and $H|I|<\int_I|f|^p$ yield
\[
 \sum_{I\in\mathcal I_z}|I|
 \le H^{-1}\sum_{I\in\mathcal I_z}\int_I|f(y,z)|^p\dd y
 \le H^{-1}\int_\R|f(y,z)|^p\dd y.
\]
Finally,
\[
 \|g\|_P^P\le\|g\|_\infty^{P-p}\|g\|_p^p
 \le(2H)^{P/p-1}\|f\|_p^p,
\]
which is the stated good-part norm after taking the $P$th root.
For any enumeration $\mathcal D_1=\{I_1,I_2,\ldots\}$, the sums
$\sum_{l\le N}b_{I_l}$ converge pointwise almost everywhere and
in $L^p$ to $b$: at most one term is nonzero at each point, and
$|b|^p\le2^{p-1}(|f|^p+|g|^p)$ is integrable.
\end{proof}

\begin{lemma}[Tails of disjoint intervals]\label{lem:interval_tails}
For pairwise disjoint bounded intervals of positive length with
$\sum_I|I|<\infty$,
put $h(x)=\sum_I(1+|x-c_I|/|I|)^{-2}$. For $1\le p<\infty$,
\[
 \|h\|_p\le C_p\left(\sum_I|I|\right)^{1/p}.
\]
\end{lemma}
\begin{proof}
For $p=1$, integrate each summand, whose integral is $2|I|$.
For $p>1$, let $v\ge0$ in $L^{p'}$. Annular decomposition and
comparison of centered intervals show, for every $y\in I$,
\[
 \int(1+|x-c_I|/|I|)^{-2}v(x)\dd x\le C|I|Mv(y).
\]
Average this inequality over $y\in I$ and sum over the disjoint
intervals. The result is at most
$C\int_{\cup I}Mv\le C_p|\cup I|^{1/p}\|v\|_{p'}$.
Duality proves the estimate for finite sums; monotone convergence
gives countable sums. In the fiber setting, apply this estimate on
each fiber and integrate its $p$th power in the remaining variables.
\end{proof}

\begin{definition}[Weak norm]\label{def:weak_norm}
For $1\le R<\infty$ define
$\|v\|_{R,\infty}=\sup_{\tau>0}\tau|\{|v|>\tau\}|^{1/R}$,
initially as an extended nonnegative number. Finite bounds below
assert the corresponding weak-$L^R$ membership.
\end{definition}

\begin{lemma}[One-fiber weak extension]\label{lem:one_fiber}
Suppose $1<P_1,P_2,P_3,R_0<\infty$,
$R_0^{-1}=\sum_jP_j^{-1}$, and
\[
 \|U(f_1,f_2,f_3)\|_{R_0}\le C_0U(u)^{100}\prod_j\|f_j\|_{P_j}.
\]
Fix $m$ and $1\le p\le P_m$ such that
$R^{-1}=p^{-1}+\sum_{j\ne m}P_j^{-1}\le1$.
These conditions imply $1\le R\le R_0$.
Then
\begin{equation}\label{eq:one_fiber}
 \|U(f_1,f_2,f_3)\|_{R,\infty}
 \le C U(u)^{100}\|f_m\|_p\prod_{j\ne m}\|f_j\|_{P_j}.
\end{equation}
The new constant depends on $C_0$ and the exponents, but is uniform
in the truncation and coefficient.
\end{lemma}
\begin{proof}
Normalize all input norms to one, with zero inputs treated separately.
Thus $\|f_m\|_p=1$ and $\|f_j\|_{P_j}=1$ for $j\ne m$.
Choose $\mathcal A_u=C_1U(u)^{100}$ with $C_1$ large enough for
the starting estimate and the fixed kernel constants, and for
$\tau>0$ put $H=(\tau/\mathcal A_u)^R$.
Apply Lemma~\ref{lem:fiber_cz} to $f_m$ at height $H$.
The good part satisfies $\|g_m\|_{P_m}\le C H^{1/p-1/P_m}$.
The starting estimate applies to this possibly nonsmooth $g_m$:
choose Schwartz approximants converging in $L^{P_m}$, and use the
preliminary truncated-operator bound to identify their output limit
with the integral defining $U$. The remaining inputs can be
approximated in their $L^{P_j}$ norms in the same way.
Write $f^{\mathrm g}=(f_1,\ldots,g_m,\ldots,f_3)$ and
$f^{\mathrm b}=(f_1,\ldots,b_m,\ldots,f_3)$ for the tuples with
only slot $m$ changed. Multilinearity gives $U(f)=U(f^{\mathrm g})+
U(f^{\mathrm b})$ almost everywhere. Chebyshev and the starting
estimate give
\[
 |\{|U(f^{\mathrm g})|>\tau/2\}|
 \le C\tau^{-R_0}\mathcal A_u^{R_0}H^{R_0(1/p-1/P_m)}
 = C\mathcal A_u^R\tau^{-R}.
\]
Indeed, subtracting the starting and target scaling identities gives
\[
 \frac1p-\frac1{P_m}=\frac1R-\frac1{R_0},\qquad
 RR_0\left(\frac1p-\frac1{P_m}\right)=R_0-R.
\]
For $I\in\mathcal D_1$, put
$2I=[c_I-|I|,c_I+|I|)$ and define
\begin{equation}\label{eq:cz_exceptional_set}
 E=\bigcup_{I\in\mathcal D_1}(2I\times S_I)
  =\{(x_m,z):x_m\in\bigcup_{I\in\mathcal I_z}2I\}.
\end{equation}
This is measurable, and Tonelli and the selected-length estimate yield
\begin{align*}
 |E|&\le\int_{E_2}\sum_{I\in\mathcal I_z}2|I|\dd z\\
 &\le2H^{-1}\int_{E_2}\int_\R|f_m(y,z)|^p\dd y\dd z
  =2H^{-1}=2\mathcal A_u^R\tau^{-R}.
\end{align*}

Outside $E$, write $d=|x_m-c_I|\ge|I|$. Cancellation of $b_I$
is the exact identity, with $s=t^{\alpha_m}$,
\begin{align*}
 (b_I*_m(\kappa_m)_{s,su_m})(x)
 =\int_I b_I(y,z)\big[&(\kappa_m)_{s,su_m}(x_m-y)\\
                      &-(\kappa_m)_{s,su_m}(x_m-c_I)\big]\dd y.
\end{align*}
For $y\in I$ and $0\le\theta\le1$,
\[
 |x_m-(c_I+\theta(y-c_I))|
 \ge d-|I|/2\ge d/2,
 \qquad |y-c_I|\le|I|/2.
\]
The mean value theorem and Lemma~\ref{lem:fiber_kernel} thus give
\begin{align*}
 &|(\kappa_m)_{s,su_m}(x_m-y)
             -(\kappa_m)_{s,su_m}(x_m-c_I)|\\
 &\hspace{15mm}\le C U(u)^{10}|I|s^{-2}(1+s^{-1}d)^{-10}.
\end{align*}
After multiplying by $\|b_I(\cdot,z)\|_1\le C H^{1/p}|I|$,
one obtains
\[
 |b_I*_m(\kappa_m)_{t^{\alpha_m},t^{\alpha_m}u_m}(x)|
 \le C U(u)^{10}H^{1/p}|I|^2t^{-2\alpha_m}
                         (1+t^{-\alpha_m}d)^{-10}.
\]
For the scale integral substitute $v=t^{-\alpha_m}d$,
so $\dd t/t=-\dd v/(\alpha_m v)$, to get
\[
 \int_0^\infty t^{-2\alpha_m}(1+t^{-\alpha_m}d)^{-10}
                  \frac{\dd t}{t}
 =\frac{d^{-2}}{\alpha_m}\int_0^\infty v(1+v)^{-10}\dd v
 =\frac{d^{-2}}{72\alpha_m}.
\]
Here $\int_0^\infty v(1+v)^{-10}\dd v=1/8-1/9=1/72$.
Also $d/|I|\ge1$ implies
\[
 \frac{|I|^2}{d^2}
 \le4\left(1+\frac d{|I|}\right)^{-2}.
\]
Bound the other two convolutions by Lemma~\ref{lem:fiber_kernel}
and sum the intervals. One obtains, almost everywhere outside $E$,
\begin{equation}\label{eq:cz_bad_pointwise}
 |U(f^{\mathrm b})(x)|
 \le C U(u)^{30}H^{1/p}h_m(x)\prod_{j\ne m}M^{(j)}f_j(x),
\end{equation}
where the measurable tail function is
\[
 h_m(x_m,z)=\sum_{I\in\mathcal D_1}\one_{S_I}(z)
       \left(1+\frac{|x_m-c_I|}{|I|}\right)^{-2}.
\]
Lemma~\ref{lem:interval_tails} on each fiber, followed by Tonelli,
gives
\begin{align*}
 \|h_m\|_p^p
 &=\int_{E_2}\|h_m(\cdot,z)\|_{L^p(\R)}^p\dd z\\
 &\le C_p\int_{E_2}\sum_{I\in\mathcal I_z}|I|\dd z
 \le C_pH^{-1}\|f_m\|_p^p=C_pH^{-1}.
\end{align*}
The maximal theorem gives
$\|M^{(j)}f_j\|_{P_j}\le C_{P_j}$ for $j\ne m$.
H\"older with $R^{-1}=p^{-1}+\sum_{j\ne m}P_j^{-1}$ therefore
implies
\[
 \left\|h_m\prod_{j\ne m}M^{(j)}f_j\right\|_R
 \le\|h_m\|_p\prod_{j\ne m}\|M^{(j)}f_j\|_{P_j}
 \le C H^{-1/p}.
\]
By \eqref{eq:cz_bad_pointwise} and Chebyshev,
\begin{align*}
 |\{x\notin E:|U(f^{\mathrm b})(x)|>\tau/2\}|
 &\le C\tau^{-R}U(u)^{30R}H^{R/p}
             \left\|h_m\prod_{j\ne m}M^{(j)}f_j\right\|_R^R\\
 &\le C\tau^{-R}U(u)^{30R}
 \le C\mathcal A_u^R\tau^{-R}.
\end{align*}
Finally, the inclusion
\[
 \{|U(f)|>\tau\}\subseteq
 \{|U(f^{\mathrm g})|>\tau/2\}\cup E\cup
 \{x\notin E:|U(f^{\mathrm b})(x)|>\tau/2\}
\]
combines the three bounds into
$\tau|\{|U(f)|>\tau\}|^{1/R}\le C\mathcal A_u$.
For countably many bad intervals, pass from finite sums using
$L^p$ convergence and the preliminary truncated-operator bound;
take a subsequence of the outputs converging almost everywhere, while
the positive tail majorants pass by Tonelli. Taking the supremum over
$\tau>0$ and restoring the input norms proves \eqref{eq:one_fiber}.
\end{proof}

\begin{exttheorem}[Multilinear Marcinkiewicz interpolation]
\label{ext:interpolation}
Let $T$ be a trilinear operator on complex simple functions of finite
measure support. Suppose
at four affinely independent input reciprocal vectors
$v^{(a)}=(1/p_{a,1},1/p_{a,2},1/p_{a,3})$, $a=0,1,2,3$,
it has weak bounds
\[
 \|T(f)\|_{r_a,\infty}\le A_a\prod_{j=1}^3\|f_j\|_{p_{a,j}},
 \quad 1<p_{a,j}<\infty,\quad1\le r_a<\infty,\quad
 r_a^{-1}=\sum_jp_{a,j}^{-1}.
\]
Let $\vartheta_a>0$, $\sum_a\vartheta_a=1$, and define
$p_j^{-1}=\sum_a\vartheta_ap_{a,j}^{-1}$,
$r^{-1}=\sum_a\vartheta_ar_a^{-1}$. If $1<r<\infty$, then
\[
 \|T(f)\|_r\le C\prod_aA_a^{\vartheta_a}\prod_j\|f_j\|_{p_j}.
\]
The constant depends only on the exponent vectors and weights.
The map extends to these $L^p$ spaces when initially defined on
their dense simple-function subspaces.
\end{exttheorem}

\begin{lemma}[Exponent simplex]\label{lem:exponent_simplex}
Suppose $0<\beta_j<1$, $\sum_{j=0}^3\beta_j=1$, and
\begin{equation}\label{eq:extended_region}
 \beta_0<\tfrac14,\quad \beta_3>\tfrac14,\quad
 \beta_1+\beta_2>\tfrac14,\quad
 \beta_1+\beta_3>\tfrac12,\quad\beta_2+\beta_3>\tfrac12.
\end{equation}
There exist $b_1,b_2,b_3>0$ with
\[
 b_j<\beta_j,\quad b_1,b_2<\tfrac14,\quad b_3<\tfrac12,
 \qquad \tfrac34<b_1+b_2+b_3<1.
\]
Put $b_0=1-\sum_{j=1}^3b_j$. Then $0<\beta_0<b_0<1/4$,
and the input vector $(\beta_1,\beta_2,\beta_3)$ is a strict convex
combination of $b=(b_1,b_2,b_3)$ and the three vectors $b+b_0e_m$.
The weights are
\begin{equation}\label{eq:weights}
 \vartheta_0=\beta_0/b_0,\qquad
 \vartheta_m=(\beta_m-b_m)/b_0\quad(m=1,2,3).
\end{equation}
\end{lemma}
\begin{proof}
Let $c=(1/4,1/4,1/2)$ and define
\begin{equation}\label{eq:capped_sum}
 S=\min(\beta_1,1/4)+\min(\beta_2,1/4)+\min(\beta_3,1/2).
\end{equation}
We prove $S>3/4$ by cases. If
$\beta_3\ge1/2$ and $\beta_1,\beta_2<1/4$, then
$S=\beta_1+\beta_2+1/2>3/4$. If $\beta_3\ge1/2$ and at
least one of $\beta_1,\beta_2$ is at least $1/4$, its capped
contribution is $1/4$, the other is strictly positive, and again
$S>3/4$. If $\beta_3<1/2$, the four remaining possibilities are
\[
 \begin{array}{c|c|c}
 &S&\text{reason }S>3/4\\ \hline
 \beta_1<1/4,\ \beta_2<1/4
    &1-\beta_0&\beta_0<1/4\\
 \beta_1\ge1/4,\ \beta_2<1/4
    &1/4+\beta_2+\beta_3&\beta_2+\beta_3>1/2\\
 \beta_1<1/4,\ \beta_2\ge1/4
    &\beta_1+1/4+\beta_3&\beta_1+\beta_3>1/2\\
 \beta_1\ge1/4,\ \beta_2\ge1/4
    &1/2+\beta_3&\beta_3>1/4.
 \end{array}
\]
Equivalently, the finite-minimum identity
\[
 S=\min_{I\subseteq\{1,2,3\}}
       \left(\sum_{j\in I}\beta_j+\sum_{j\notin I}c_j\right)
 \]
expands into the eight alternatives
\begin{align*}
 S=\min\{&1,\ \beta_1+3/4,\ \beta_2+3/4,\ \beta_3+1/2,\\
         &\beta_1+\beta_2+1/2,\ \beta_1+\beta_3+1/4,
            \ \beta_2+\beta_3+1/4,\ 1-\beta_0\}.
\end{align*}
Each entry exceeds $3/4$ by positivity or
\eqref{eq:extended_region}.
Make the choice constructive by setting
\[
 \epsilon=\frac12\min\left\{
                  \min_{1\le j\le3}\min(\beta_j,c_j),
                              \frac{S-3/4}{3}\right\}>0,
 \qquad b_j=\min(\beta_j,c_j)-\epsilon.
\]
Then $b_j>0$, $b_j<\min(\beta_j,c_j)$, and
\[
 \frac34<S-3\epsilon=\sum_{j=1}^3b_j
           <\sum_{j=1}^3\beta_j=1-\beta_0<1.
\]
With $b_0=1-\sum_{j=1}^3b_j$ this gives
$0<\beta_0<b_0<1/4$. It also gives
$b_3=1-b_0-b_1-b_2>1/4$ since the three subtracted terms are
each less than $1/4$.
Set $\delta_m=\beta_m-b_m>0$. Then
\[
 \sum_{m=1}^3\delta_m
 =(1-\beta_0)-(1-b_0)=b_0-\beta_0,
 \qquad\sum_{a=0}^3\vartheta_a
 =\frac{\beta_0+\sum_m\delta_m}{b_0}=1.
\]
For each input coordinate $j$, the strict convex-combination identity is
\[
 \vartheta_0b_j+\sum_{m=1}^3\vartheta_m
                  (b_j+b_0\one_{\{m=j\}})
 =b_j+b_0\vartheta_j=\beta_j.
\]
The matching output reciprocal is
\[
 \vartheta_0(1-b_0)+\sum_{m=1}^3\vartheta_m
 =1-b_0\vartheta_0=1-\beta_0.
\]
The simplex has independent edge vectors $b_0e_1,b_0e_2,b_0e_3$.
\end{proof}

\begin{theorem}[Extended model range]\label{thm:extended_model}
Let $1<p_j<\infty$ with $\sum_jp_j^{-1}=1$, and suppose
$\beta_j=p_j^{-1}$ satisfies \eqref{eq:extended_region}. Then
\[
 |\Lambda_{u,c}(f)|\le C_{\alpha,p}U(u)^{100}\prod_{j=0}^3\|f_j\|_{p_j}
\]
for complex Schwartz inputs, uniformly in $u,c$. The same estimate
holds for all scale truncations, and defines a unique continuous form
on the product of the four $L^{p_j}$ spaces.
\end{theorem}
\begin{proof}
Choose $b$ by Lemma~\ref{lem:exponent_simplex}. Put
$P_j=1/b_j$ for $j=1,2,3$ and $R_0=1/(1-b_0)$.
Theorem~\ref{thm:initial_model}, with zeroth exponent $1/b_0$,
gives by $L^p$ duality
\[
 \|U^{a,b}_{u,c}(f)\|_{R_0}\le C U(u)^{100}\prod_j\|f_j\|_{P_j}.
\]
For complex functions the bilinear pairing $\int f_0U$ gives the
same dual norm, by replacing the usual test function by its conjugate.
The hypotheses of the initial range hold since $b_0,b_1,b_2<1/4$
and $b_3<1/2$; also $b_3>1/4$ follows by subtracting the first
three coordinates from $1$.
For each $m\in\{1,2,3\}$ define the complete endpoint tuple by
\[
 P_j^{(m)}=\begin{cases}
        (b_m+b_0)^{-1},&j=m,\\
        P_j,&j\ne m.
        \end{cases}
\]
The identities
\[
 b_m<b_m+b_0=1-\sum_{j\ne m}b_j<1,\qquad
 \sum_{j=1}^3\frac1{P_j^{(m)}}=
           b_0+\sum_{j=1}^3b_j=1
\]
show $1<P_m^{(m)}<P_m$ and output exponent $1$. Apply
Lemma~\ref{lem:one_fiber} with $p=P_m^{(m)}$ to obtain
\begin{equation}\label{eq:interpolation_weak_endpoints}
 \|U^{a,b}_{u,c}(f_1,f_2,f_3)\|_{1,\infty}
 \le C_m U(u)^{100}\prod_{j=1}^3\|f_j\|_{P_j^{(m)}}
                  \qquad(m=1,2,3).
\end{equation}
These are the three vertices $b+b_0e_m$.
The base strong bound is also a weak bound at $b$.
External Theorem~\ref{ext:interpolation}, with the weights
\eqref{eq:weights}, gives the strong estimate at the target inputs
and output $R$ with $R^{-1}=1-\beta_0$:
\begin{equation}\label{eq:extended_operator_bound}
 \|U^{a,b}_{u,c}(f_1,f_2,f_3)\|_R
 \le C_{\alpha,p} U(u)^{100}\prod_{j=1}^3\|f_j\|_{p_j}.
\end{equation}
Indeed all endpoint constants have factor $U(u)^{100}$, and
\[
 \prod_{a=0}^3(U(u)^{100})^{\vartheta_a}
 =U(u)^{100\sum_a\vartheta_a}=U(u)^{100}.
\]
Since $R^{-1}+p_0^{-1}=1$, H\"older gives
\begin{align*}
 |\Lambda^{a,b}_{u,c}(f)|
 &=\left|\int_{E_3}f_0(x)U^{a,b}_{u,c}(f_1,f_2,f_3)(x)\dd x\right|\\
 &\le\|f_0\|_{p_0}\|U^{a,b}_{u,c}(f_1,f_2,f_3)\|_R\\
 &\le C_{\alpha,p}U(u)^{100}\prod_{j=0}^3\|f_j\|_{p_j}.
\end{align*}
The constants are independent of $a,b$; take the scale limit for
Schwartz inputs by Lemma~\ref{lem:model_convergence}, then use density.
\end{proof}

\section{Cone decomposition and the multiplier theorem}

\begin{definition}[Models in each active coordinate]\label{def:permuted_model}
For $i,j\in\{1,2,3\}$ define
\[
 \kappa_j^{(i)}=
 \begin{cases}
  \h*\h*\psi,&j=i,\\
  \g*\varphi,&j\ne i.
 \end{cases}
\]
The corresponding truncated and full-scale forms are
\begin{align*}
 \Lambda^{(i),a,b}_{u,c}(f)
 &=\int_a^b c(t)\int_{E_3}f_0(x)
       \prod_{j=1}^3
       \big(f_j*_j(\kappa_j^{(i)})_{t^{\alpha_j},t^{\alpha_j}u_j}\big)(x)
                                      \dd x\frac{\dd t}{t},\\
 \Lambda^{(i)}_{u,c}(f)
 &=\lim_{a\downarrow0,\ b\uparrow\infty}\Lambda^{(i),a,b}_{u,c}(f).
\end{align*}
The last limit is initially taken on Schwartz tuples.
For $i=3$ these are the forms of Definition~\ref{def:model}.
Translation commutes with convolution, so, for example,
\[
 (\kappa_i^{(i)})_{t^{\alpha_i},t^{\alpha_i}u_i}
      =(\h*\h*\psi_{1,u_i})_{(t^{\alpha_i})}.
\]
\end{definition}

\begin{lemma}[Coordinate permutation]\label{lem:permutation}
Choose a permutation $\sigma$ with $\sigma(3)=i$ and put
\[
 P_\sigma x=(x_{\sigma(1)},x_{\sigma(2)},x_{\sigma(3)}),\qquad
 \widetilde\alpha_j=\alpha_{\sigma(j)},\quad
 \widetilde u_j=u_{\sigma(j)},
\]
\[
 \widetilde f_0=f_0\circ P_\sigma^{-1},\qquad
 \widetilde f_j=f_{\sigma(j)}\circ P_\sigma^{-1}\quad(1\le j\le3).
\]
Writing the anisotropy as an extra subscript, one has
\[
 \Lambda^{(i),a,b}_{\alpha,u,c}(f)
      =\Lambda^{(3),a,b}_{\widetilde\alpha,\widetilde u,c}(\widetilde f),
 \qquad
 \Lambda^{(i)}_{\alpha,u,c}(f)
      =\Lambda^{(3)}_{\widetilde\alpha,\widetilde u,c}(\widetilde f).
\]
Lebesgue measure and Schwartz regularity are preserved, as are the
corresponding $L^p$ norms and $U(\widetilde u)=U(u)$.
\end{lemma}
\begin{proof}
The identities
\[
 P_\sigma e_{\sigma(j)}=e_j,\qquad
 P_\sigma D_t^\alpha=D_t^{\widetilde\alpha}P_\sigma,
 \qquad |\det P_\sigma|=1
\]
give the result by setting $z=P_\sigma x$ and relabeling each fiber
integration variable. The $\sigma(j)$th factor becomes
\[
 \int_\R\widetilde f_j(z-ye_j)
     (\kappa_j^{(3)})_{t^{\widetilde\alpha_j},
                                 t^{\widetilde\alpha_j}\widetilde u_j}(y)\dd y.
\]
The coefficient $c(t)$ is unchanged. The change of variables has
Jacobian one in absolute value, and $P_\sigma$ is an isometry, which
proves the norm and regularity assertions. The full-scale limits
exist by Lemma~\ref{lem:model_convergence} for the permuted data;
passing to those limits gives the second identity.
\end{proof}

\begin{definition}[Auxiliary cutoff and localized symbol]\label{def:localized_symbol}
Use $b=\widehat\B$ from Definition~\ref{def:bumps}, and put
$\zeta=b/\int b$. This is a nonnegative even smooth probability
density supported in $[-1,1]$. Define
\[
 \widetilde\Phi=\one_{[-5/2,5/2]}*\zeta_{(1/2)},\qquad
 \widetilde\Psi=(\one_{[3/4,9/4]}+\one_{[-9/4,-3/4]})*\zeta_{(1/4)}.
\]
Then $\widetilde\Phi=1$ on $[-2,2]$, with support in $[-3,3]$,
and $\widetilde\Psi=1$ on $[-2,-1]\cup[1,2]$, with support in
$[-5/2,-1/2]\cup[1/2,5/2]$.
For $i=1,2,3$, $t>0$ define
\[
 \chi_i(\eta)=\widetilde\Psi(\eta_i)
          \prod_{j\ne i}\widetilde\Phi(\eta_j)\g(\eta_j)^{-1},
 \qquad m_{i,t}(\eta)=m(D_t^{-1}\eta)\chi_i(\eta).
\]
Set $m_{i,t}=0$ near the origin, where the cutoff vanishes.
\end{definition}

\begin{lemma}[Calder\'on identity and three cones]\label{lem:calderon}
For $\xi_i\ne0$,
\[
 \int_0^\infty\alpha_i\Psi(t^{\alpha_i}\xi_i)
                 \h(t^{\alpha_i}\xi_i)^2\frac{\dd t}{t}=c_\Psi.
\]
For every $\xi\ne0$, putting $\eta=D_t\xi$ gives
\begin{equation}\label{eq:cone_identity}
 m(\xi)=c_\Psi^{-3}\sum_{i=1}^3\alpha_i
     \int_0^\infty m(\xi)\Psi(\eta_i)\h(\eta_i)^2
                          \prod_{j\ne i}\Phi(\eta_j)\frac{\dd t}{t}.
\end{equation}
\end{lemma}
\begin{proof}
For the first formula substitute $v=t^{\alpha_i}|\xi_i|$ and use
evenness: $\dd v/v=\alpha_i\dd t/t$.
If $\xi_1\xi_2\xi_3\ne0$, multiply the three one-dimensional
identities to obtain
\[
 m(\xi)=c_\Psi^{-3}\int_{(0,\infty)^3}m(\xi)
      \prod_{j=1}^3\alpha_j\Psi(t_j^{\alpha_j}\xi_j)
                                  \h(t_j^{\alpha_j}\xi_j)^2
                          \prod_{j=1}^3\frac{\dd t_j}{t_j}.
\]
This integral is absolutely convergent, since its absolute integral
is $|m(\xi)|c_\Psi^3$ before multiplication by $c_\Psi^{-3}$.
Define the measurable regions
\[
 \mathcal C_i=\{(t_1,t_2,t_3)\in(0,\infty)^3:
                      t_i<t_j\text{ for every }j\ne i\}.
\]
They are disjoint and cover the positive octant except for the null
set contained in $\bigcup_{j<k}\{t_j=t_k\}$.
For fixed $t_i=t$, integrate the other two scales over $(t,\infty)$.
The substitution $v=s^{\alpha_j}|\xi_j|$ gives
\[
 \int_t^\infty\alpha_j\Psi(s^{\alpha_j}\xi_j)
           \h(s^{\alpha_j}\xi_j)^2\frac{\dd s}{s}
       =\int_{t^{\alpha_j}|\xi_j|}^\infty
                \Psi(v)\h(v)^2\frac{\dd v}{v}
       =\Phi(t^{\alpha_j}\xi_j).
\]
Summing over $i$ yields \eqref{eq:cone_identity} on this set.

For a direct verification valid also on coordinate planes,
differentiate the product $\prod_{j=1}^3\Phi(t^{\alpha_j}\xi_j)$.
Each nonzero coordinate satisfies
\[
 -t\partial_t\Phi(t^{\alpha_j}\xi_j)
       =\alpha_j\Psi(t^{\alpha_j}\xi_j)\h(t^{\alpha_j}\xi_j)^2;
\]
for a zero coordinate both sides vanish. The product is $c_\Psi^3$
for all sufficiently small $t$ and is zero for all sufficiently
large $t$, since at least one coordinate of $\xi$ is nonzero.
Integrate its derivative and multiply by $c_\Psi^{-3}m(\xi)$.
Every integrand is supported in a compact positive scale interval
for each fixed nonzero $\xi$.
The scale restriction is uniform on compact subsets away from zero:
for such a set $K$, put
\[
 M_K=\max(1,\sup_{\xi\in K}|\xi|),\qquad
 \delta_K=\inf_{\xi\in K}\max_j|\xi_j|>0,
\]
\[
 a_K=\min_j M_K^{-1/\alpha_j},\qquad
 b_K=\max_j(2/\delta_K)^{1/\alpha_j}.
\]
On a nonzero integrand, the active factor forces
$t^{\alpha_i}|\xi_i|\ge1$, hence $t\ge a_K$.
All three frequency factors force $t^{\alpha_j}|\xi_j|\le2$;
choosing a coordinate with $|\xi_j|\ge\delta_K$ gives $t\le b_K$.
This also justifies differentiation of each cone integral in $\xi$
and proves its smoothness off the origin.
\end{proof}

\begin{lemma}[Uniform localized derivatives]\label{lem:symbol_derivatives}
The symbol $m_{i,t}$ is smooth, supported in $[-3,3]^3$ and away
from $\eta_i=0$, and for $|\beta|\le110$,
\[
 \sup_{i,t>0}\|\partial^\beta m_{i,t}\|_\infty\le C_{\alpha,\beta}M.
\]
It extends to a smooth function periodic with period $8$ in each
coordinate by periodization, without overlap near the boundary of
the fundamental cube $[-4,4]^3$.
\end{lemma}
\begin{proof}
Smoothness and plateau properties of the auxiliary cutoffs follow by
differentiating their convolutions with the compactly supported smooth
density $\zeta$. Their supports have the stated margins.
For $\eta\ne0$, the Leibniz rule and the coordinate chain rule give
the full expansion
\[
 \partial_\eta^\beta m_{i,t}(\eta)
 =\sum_{\gamma\le\beta}\binom{\beta}{\gamma}
 t^{-\alpha\cdot\gamma}(\partial^\gamma m)(D_t^{-1}\eta)
                  \partial^{\beta-\gamma}\chi_i(\eta).
\]
In a neighborhood of $\eta=0$ the cutoff and the zero-extended
symbol vanish identically, so every derivative there is zero;
no derivative of $m$ at the origin is used.
On the support of the cutoff derivatives, $|\eta_i|\ge1/2$ and
$\rho_\alpha(\eta)\ge(1/2)^{1/\alpha_i}$. The multiplier hypothesis
and $\rho_\alpha(D_t^{-1}\eta)=t^{-1}\rho_\alpha(\eta)$ give
\[
 t^{-\alpha\cdot\gamma}|(\partial^\gamma m)(D_t^{-1}\eta)|
 \le M\rho_\alpha(\eta)^{-\alpha\cdot\gamma}.
\]
All cutoff derivatives are bounded on the fixed compact support;
division by the positive Gaussian causes no singularity there.
Sum the finitely many terms to obtain the bound.
The periodic extension is the locally finite sum
\[
 m_{i,t}^{\mathrm{per}}(\eta)
       =\sum_{k\in\Z^3}m_{i,t}(\eta+8k).
\]
On $[-4,4]^3$ only the term $k=0$ can be nonzero. The compact
support leaves a neighborhood of every cell boundary on which
the adjacent terms vanish, so all derivatives match there.
\end{proof}

\begin{definition}[Fourier coefficients]\label{def:fourier_coefficients}
For $\nu\in\Z^3$, put
\[
 a_i(t,\nu)=8^{-3}\int_{[-4,4]^3}m_{i,t}(\eta)
                                   e^{-2\pi i\nu\cdot\eta/8}\dd\eta.
\]
\end{definition}

\begin{lemma}[Fourier-series expansion and summability]\label{lem:fourier_series}
There is $C_\alpha\ge1$ such that
\begin{equation}\label{eq:coefficient_decay}
 |a_i(t,\nu)|\le C_\alpha M(1+|\nu|)^{-110}
\end{equation}
for all $i,t,\nu$. Each coefficient is measurable in $t$ and, on
$[-4,4]^3$, the uniformly absolutely convergent Fourier series gives
\[
 m_{i,t}(\eta)=\sum_{\nu\in\Z^3}a_i(t,\nu)e^{2\pi i\nu\cdot\eta/8}.
\]
For every $s>3$, $\sum_{\nu\in\Z^3}(1+|\nu|)^{-s}<\infty$.
\end{lemma}
\begin{proof}
Integrate by parts using $(1-\Delta_\eta)^{55}$. Boundary terms
vanish by the support margin, and
\[
 (1+\pi^2|\nu|^2/16)^{55}|a_i(t,\nu)|
 \le8^{-3}\|(1-\Delta)^{55}m_{i,t}\|_{L^1([-4,4]^3)}.
\]
Lemma~\ref{lem:symbol_derivatives} bounds the right side by $CM$;
this uses exactly $110$ derivatives. Measurability follows from
measurability of the joint integrand and parameter integration.
Lattice points in a ball of radius $2^{k+1}$ number at most
$C2^{3k}$. Splitting into dyadic shells proves summability for $s>3$.
Thus the Fourier series converges absolutely and uniformly. It has
the same Fourier coefficients as the periodic symbol. Uniqueness
of Fourier coefficients, for example by convergence of the product
Fej\'er means of continuous periodic functions, identifies its sum
with the symbol.
\end{proof}

\begin{theorem}[Cone decomposition of the form]\label{thm:cone}
If $M>0$, choose $A_0=C_\alpha M>0$ large enough for
\eqref{eq:coefficient_decay}, and put
\[
 w_i(\nu)=c_\Psi^{-3}\alpha_iA_0(1+|\nu|)^{-110},\qquad
 c_{i,\nu}(t)=-\frac{a_i(t,\nu)}{A_0(1+|\nu|)^{-110}}.
\]
Then the coefficients $c_{i,\nu}$ are measurable with modulus at
most one, and for Schwartz tuples
\begin{equation}\label{eq:form_decomposition}
 \Lambda_m(f)=\sum_{i=1}^3\sum_{\nu\in\Z^3}
                  w_i(\nu)\Lambda^{(i)}_{\nu/8,c_{i,\nu}}(f).
\end{equation}
The expansion is absolutely convergent. If $M=0$, the form is zero.
\end{theorem}
\begin{proof}
On the support of the $i$th integrand of \eqref{eq:cone_identity},
the auxiliary cutoffs are one. Hence
\[
 m(\xi)\Psi(\eta_i)\h(\eta_i)^2\prod_{j\ne i}\Phi(\eta_j)
 =m_{i,t}(\eta)\Psi(\eta_i)\h(\eta_i)^2
                           \prod_{j\ne i}\g(\eta_j)\Phi(\eta_j).
\]
Insert the Fourier series of $m_{i,t}$. To verify the signs, use
frequency variables $\zeta^{(j)}\in E_3$ for $f_j$, $j=1,2,3$,
and set $\xi_j=-\zeta^{(j)}_j$.
By \eqref{eq:frequency_form}, integration against $f_0$ produces
$\widehat f_0(-\zeta^{(1)}-\zeta^{(2)}-\zeta^{(3)})$.
The exceptional set $\xi=0$ is a null coordinate subspace of $E_9$,
so the punctured cone identity applies almost everywhere in this integral.
The two kinds of fiber kernel multipliers, evaluated at $-\eta_j$,
are
\begin{align*}
 \widehat{\g*\varphi_{1,\nu_j/8}}(-\eta_j)
    &=\g(\eta_j)\Phi(\eta_j)e^{2\pi i\nu_j\eta_j/8},\quad j\ne i,\\
 \widehat{\h*\h*\psi_{1,\nu_i/8}}(-\eta_i)
    &=-\h(\eta_i)^2\Psi(\eta_i)e^{2\pi i\nu_i\eta_i/8}.
\end{align*}
The second minus sign follows from $\widehat\h=-i\h$ and is
absorbed into $c_{i,\nu}$. The remaining constants are exactly
$w_i(\nu)$.

For absolute convergence, the nonnegative scale integral
\[
 \int_0^\infty\Psi(t^{\alpha_i}\xi_i)\h(t^{\alpha_i}\xi_i)^2
           \prod_{j\ne i}\g(t^{\alpha_j}\xi_j)\Phi(t^{\alpha_j}\xi_j)
                                                \frac{\dd t}{t}
\]
is at most $c_\Psi^3/\alpha_i$ when $\xi_i\ne0$ and is zero when
$\xi_i=0$. This uses $0\le\Phi\le c_\Psi$ and $0<\g\le1$.
The product
\[
 |\widehat f_0(-\zeta^{(1)}-\zeta^{(2)}-\zeta^{(3)})|
                         \prod_{j=1}^3|\widehat f_j(\zeta^{(j)})|
\]
is integrable in $E_9$, bounded in integral by
$\|\widehat f_0\|_\infty\prod_{j=1}^3\|\widehat f_j\|_1$.
Together with \eqref{eq:coefficient_decay} and lattice summability,
these bounds justify Fubini between frequencies, scales, and modes.
Apply Fourier inversion first with finitely many modes and compact
scale intervals, then pass to the limit using these majorants.
Lemma~\ref{lem:model_convergence} identifies each scale limit with
its model form. This proves \eqref{eq:form_decomposition}.
If $M=0$, the zeroth derivative bound makes $m=0$ almost everywhere.
\end{proof}

\begin{proof}[Proof of Theorem~\ref{thm:main}]
Let $\beta_j=1/p_j$. The hypotheses give
\[
 \beta_0<\frac14,\qquad
 \beta_1,\beta_2,\beta_3>\frac14,\qquad
 \sum_{j=0}^3\beta_j=1.
\]
Fix an active coordinate $i$ and write $\{i,j,k\}=\{1,2,3\}$.
The inequalities required by the extended model estimate are
\[
 \beta_i>\frac14,\qquad
 \beta_i+\beta_j>\frac12,\qquad
 \beta_i+\beta_k>\frac12,\qquad
 \beta_j+\beta_k>\frac12>\frac14.
\]
Together with the condition on $\beta_0$ and the reciprocal sum,
these are \eqref{eq:extended_region} after permutation.
Lemmas~\ref{lem:permutation} and Theorem~\ref{thm:extended_model}
therefore give
\[
 |\Lambda^{(i)}_{\nu/8,c_{i,\nu}}(f)|
       \le C_{\alpha,p}U(\nu/8)^{100}\prod_{j=0}^3\|f_j\|_{p_j}.
\]
Since $U(\nu/8)\le C(1+|\nu|)$ and
$w_i(\nu)\le C_\alpha M(1+|\nu|)^{-110}$, each term in
\eqref{eq:form_decomposition} is bounded by
\[
 C_{\alpha,p}M(1+|\nu|)^{-10}\prod_{j=0}^3\|f_j\|_{p_j}.
\]
Sum using Lemma~\ref{lem:fourier_series}, with $s=10>3$.
The case $M=0$ has already been handled.
\end{proof}

\end{document}
